
\documentclass{article}
\usepackage{colortbl}
\usepackage{makecell}
\usepackage{multirow}
\usepackage{supertabular}

\begin{document}

\newcounter{utterance}

\centering \large Interaction Transcript for game `clean\_up', experiment `1\_medium\_7obj\_en', episode 1 with deepseek{-}chat{-}v3.1{-}t1.0.
\vspace{24pt}

{ \footnotesize  \setcounter{utterance}{1}
\setlength{\tabcolsep}{0pt}
\begin{supertabular}{c@{$\;$}|p{.15\linewidth}@{}p{.15\linewidth}p{.15\linewidth}p{.15\linewidth}p{.15\linewidth}p{.15\linewidth}}
    \# & $\;$A & \multicolumn{4}{c}{Game Master} & $\;\:$B\\
    \hline

    \theutterance \stepcounter{utterance}  
    & & \multicolumn{4}{p{0.6\linewidth}}{
        \cellcolor[rgb]{0.9,0.9,0.9}{
            \makecell[{{p{\linewidth}}}]{
                \texttt{\tiny{[P1$\langle$GM]}}
                \texttt{I am your game master, and you are playing a collaborative game with the following grid as game board:} \\
\\ 
\texttt{```} \\
\texttt{ 1234567} \\
\texttt{╔═══╤═╤═╗} \\
\texttt{║◌◌◌│◌│◌║ 1} \\
\texttt{║◌◌M│◌├─╢ 2} \\
\texttt{║AT◌├─┘◌║ 3} \\
\texttt{║◌◌P│◌◌◌║ 4} \\
\texttt{╟─┬─┼───╢ 5} \\
\texttt{║◌│◌│S◌O║ 6} \\
\texttt{║◌│◌│◌D◌║ 7} \\
\texttt{╚═╧═╧═══╝} \\
\\ 
\texttt{```} \\
\\ 
\texttt{* The upper edge displays x{-}coordinates increasing to the right, and the right edge y{-}coordinates increasing downward.} \\
\texttt{* The following objects are randomly placed on your grid: 'P', 'O', 'T', 'S', 'D', 'A', 'M'.} \\
\\ 
\\ 
\texttt{The other player sees a variation of the game board, where the objects are placed at different random locations. You cannot see the other player's board, and they cannot see yours.} \\
\\ 
\\ 
\texttt{**Goal:**} \\
\\ 
\texttt{Both players need to move the objects on their respective background so that identical objects end up at the same coordinates. You have to communicate with the other player to agree upon a common goal configuration.} \\
\\ 
\\ 
\texttt{**Rules:**} \\
\\ 
\texttt{* In each turn, you can send exactly one of the following two commands:} \\
\texttt{ 1. `SAY: <MESSAGE>`: to send a message (everything up until the next line break) to the other player. I will forward it to your partner.} \\
\texttt{2. `MOVE: <OBJECT>, (<X>, <Y>)`: to move an object to a new position, where `<X>` is the column and `<Y>` is the row. I will inform you if your move was successful or not.} \\
\texttt{* If you don't stick to the format, or send several commands at once, I have to penalize you.} \\
\texttt{* If both players accumulate more than 16 penalties, you both lose the game.} \\
\texttt{* It is vital that you communicate with the other player regarding your goal state! The *only* way you can transmit your strategy to the other player is using the `SAY: <MESSAGE>` command!} \\
\\ 
\\ 
\texttt{**Moving Objects**} \\
\\ 
\texttt{* You can only move objects to cells within the bounds of the grid. The target cell must be empty, i.e., it must only contain the symbol '◌'.} \\
\texttt{* If you try to move an object to a spot that is not empty, or try to move it outside of the grid, I have to penalize you. You get another try.} \\
\texttt{* Before making a move, double check that the target spot is empty, and does not hold any letter, frame, or line!} \\
\\ 
\\ 
\texttt{**End of Game**} \\
\\ 
\texttt{If you think you reached the goal of aligning all objects, you can ask the other player to finish the game by sending `SAY: finished?`. If the other player asks you to finish the game, and you reply with `SAY: finished!`, the game will end.} \\
\\ 
\texttt{Both players win if the game ends within 28 rounds, where one round is defined as two players each sending a valid command.} \\
\\ 
\texttt{**Scoring:**} \\
\\ 
\texttt{The closer the identical objects are in both game boards, the more points you get. Penalties reduce your points. Can you beat the record?} \\
\\ 
\\ 
\texttt{Please send a message to the other player to start the game!} \\
            }
        }
    }
    & \\ \\

    \theutterance \stepcounter{utterance}  
    & \multicolumn{4}{p{0.6\linewidth}}{
        \cellcolor[rgb]{0.9,0.9,0.9}{
            \makecell[{{p{\linewidth}}}]{
                \texttt{\tiny{[P1$\rangle$GM]}}
                \texttt{SAY: Hello! Let's coordinate to align our objects. I'll start by suggesting a common goal configuration. How about we place objects in a standard order? For example, let's put all objects in row 4 from left to right: A, D, M, O, P, S, T. So coordinates: A at (1,4), D at (2,4), M at (3,4), O at (4,4), P at (5,4), S at (6,4), T at (7,4). What do you think? If you agree, I'll start moving.} \\
            }
        }
    }
    & & \\ \\

    \theutterance \stepcounter{utterance}  
    & & & \multicolumn{2}{p{0.3\linewidth}}{
        \cellcolor[rgb]{0.9,0.9,0.9}{
            \makecell[{{p{\linewidth}}}]{
                \texttt{\tiny{[GM$|$GM]}}
                \texttt{SAY: Hello! Let's coordinate to align our objects. I'll start by suggesting a common goal configuration. How about we place objects in a standard order? For example, let's put all objects in row 4 from left to right: A, D, M, O, P, S, T. So coordinates: A at (1,4), D at (2,4), M at (3,4), O at (4,4), P at (5,4), S at (6,4), T at (7,4). What do you think? If you agree, I'll start moving.} \\
            }
        }
    }
    & & \\ \\

    \theutterance \stepcounter{utterance}  
    & & \multicolumn{4}{p{0.6\linewidth}}{
        \cellcolor[rgb]{0.9,0.9,0.9}{
            \makecell[{{p{\linewidth}}}]{
                \texttt{\tiny{[GM$\rangle$P2]}}
                \texttt{I am your game master, and you are playing a collaborative game with the following grid as game board:} \\
\\ 
\texttt{```} \\
\texttt{ 1234567} \\
\texttt{╔═══╤═╤═╗} \\
\texttt{║◌MS│◌│◌║ 1} \\
\texttt{║◌O◌│◌├─╢ 2} \\
\texttt{║◌◌◌├─┘◌║ 3} \\
\texttt{║◌◌◌│◌◌◌║ 4} \\
\texttt{╟─┬─┼───╢ 5} \\
\texttt{║◌│◌│◌◌D║ 6} \\
\texttt{║T│◌│◌PA║ 7} \\
\texttt{╚═╧═╧═══╝} \\
\\ 
\texttt{```} \\
\\ 
\texttt{* The upper edge displays x{-}coordinates increasing to the right, and the right edge y{-}coordinates increasing downward.} \\
\texttt{* The following objects are randomly placed on your grid: 'P', 'O', 'T', 'S', 'D', 'A', 'M'.} \\
\\ 
\\ 
\texttt{The other player sees a variation of the game board, where the objects are placed at different random locations. You cannot see the other player's board, and they cannot see yours.} \\
\\ 
\\ 
\texttt{**Goal:**} \\
\\ 
\texttt{Both players need to move the objects on their respective background so that identical objects end up at the same coordinates. You have to communicate with the other player to agree upon a common goal configuration.} \\
\\ 
\\ 
\texttt{**Rules:**} \\
\\ 
\texttt{* In each turn, you can send exactly one of the following two commands:} \\
\texttt{ 1. `SAY: <MESSAGE>`: to send a message (everything up until the next line break) to the other player. I will forward it to your partner.} \\
\texttt{2. `MOVE: <OBJECT>, (<X>, <Y>)`: to move an object to a new position, where `<X>` is the column and `<Y>` is the row. I will inform you if your move was successful or not.} \\
\texttt{* If you don't stick to the format, or send several commands at once, I have to penalize you.} \\
\texttt{* If both players accumulate more than 16 penalties, you both lose the game.} \\
\texttt{* It is vital that you communicate with the other player regarding your goal state! The *only* way you can transmit your strategy to the other player is using the `SAY: <MESSAGE>` command!} \\
\\ 
\\ 
\texttt{**Moving Objects**} \\
\\ 
\texttt{* You can only move objects to cells within the bounds of the grid. The target cell must be empty, i.e., it must only contain the symbol '◌'.} \\
\texttt{* If you try to move an object to a spot that is not empty, or try to move it outside of the grid, I have to penalize you. You get another try.} \\
\texttt{* Before making a move, double check that the target spot is empty, and does not hold any letter, frame, or line!} \\
\\ 
\\ 
\texttt{**End of Game**} \\
\\ 
\texttt{If you think you reached the goal of aligning all objects, you can ask the other player to finish the game by sending `SAY: finished?`. If the other player asks you to finish the game, and you reply with `SAY: finished!`, the game will end.} \\
\\ 
\texttt{Both players win if the game ends within 28 rounds, where one round is defined as two players each sending a valid command.} \\
\\ 
\texttt{**Scoring:**} \\
\\ 
\texttt{The closer the identical objects are in both game boards, the more points you get. Penalties reduce your points. Can you beat the record?} \\
\\ 
\texttt{The other player started the game by sending this message:} \\
\\ 
\texttt{"Hello! Let's coordinate to align our objects. I'll start by suggesting a common goal configuration. How about we place objects in a standard order? For example, let's put all objects in row 4 from left to right: A, D, M, O, P, S, T. So coordinates: A at (1,4), D at (2,4), M at (3,4), O at (4,4), P at (5,4), S at (6,4), T at (7,4). What do you think? If you agree, I'll start moving."} \\
\\ 
\texttt{What is your first command?} \\
            }
        }
    }
    & \\ \\

    \theutterance \stepcounter{utterance}  
    & & & \multicolumn{4}{p{0.6\linewidth}}{
        \cellcolor[rgb]{0.9,0.9,0.9}{
            \makecell[{{p{\linewidth}}}]{
                \texttt{\tiny{[GM$\langle$P2]}}
                \texttt{Let me analyze the situation. The other player has proposed a common goal configuration: all objects in row 4 from left to right: A at (1,4), D at (2,4), M at (3,4), O at (4,4), P at (5,4), S at (6,4), T at (7,4).} \\
\\ 
\texttt{However, looking at my grid, I can see that row 4 currently has three '◌' symbols and no objects. Specifically, row 4 is: `║◌◌◌│◌◌◌║ 4`. This means there are no objects in row 4 right now. Also, the grid has boundaries (like lines and corners) that might interfere with placements.} \\
\\ 
\texttt{But importantly, the proposed positions (1,4) to (7,4) are all within the grid and should be empty in my board (since row 4 is entirely empty of objects). So it's feasible.} \\
\\ 
\texttt{However, I should verify if all these target cells are indeed empty. Let's check:} \\
\\ 
\texttt{{-} (1,4): column 1, row 4 {-} this is the first cell in row 4, which is '◌' (empty).} \\
\texttt{{-} (2,4): also '◌'} \\
\texttt{{-} (3,4): also '◌'} \\
\texttt{{-} (4,4): this is the middle, but note there is a vertical line at column 4? Actually, looking at the grid, the columns are separated by lines. For example, in row 1: `║◌MS│◌│◌║` {-} the vertical lines are at column boundaries.} \\
\\ 
\texttt{Wait, let's understand the grid layout:} \\
\\ 
\texttt{The grid has 7 columns and 7 rows. The vertical lines are drawn between columns? Actually, the grid is drawn with box{-}drawing characters.} \\
\\ 
\texttt{Looking at row 1: `║◌MS│◌│◌║ 1`} \\
\texttt{This means:} \\
\texttt{{-} Column 1: ◌} \\
\texttt{{-} Column 2: M} \\
\texttt{{-} Column 3: S} \\
\texttt{{-} Column 4: │ (this is a line, not a cell)} \\
\texttt{Wait no, the lines are separators. The actual cells are between the lines.} \\
\\ 
\texttt{Actually, the grid is represented with the cells being the spaces between the lines. The numbers on top (1234567) indicate the column numbers.} \\
\\ 
\texttt{So for row 1:} \\
\texttt{{-} The left boundary is at x=0, but column 1 is at x=1.} \\
\texttt{{-} The string "◌MS│◌│◌" might be misleading.} \\
\\ 
\texttt{Let's parse the grid properly:} \\
\\ 
\texttt{The grid is 7x7. The top line "╔═══╤═╤═╗" indicates the top border, with openings.} \\
\\ 
\texttt{Better to think: each cell is one character wide? But some cells have multiple characters? Actually, no, each cell is one symbol.} \\
\\ 
\texttt{Actually, the grid is drawn with each cell being one character. The vertical lines are drawn between cells.} \\
\\ 
\texttt{For row 1: `║◌MS│◌│◌║` {-} this is confusing. Probably it's:} \\
\texttt{{-} First cell (col1): '◌'} \\
\texttt{{-} Second cell (col2): 'M'} \\
\texttt{{-} Third cell (col3): 'S'} \\
\texttt{{-} Then a vertical line '│' (which is a separator, not a cell)} \\
\texttt{{-} Then col4: '◌'} \\
\texttt{{-} Then a vertical line '│'} \\
\texttt{{-} Then col5: '◌'} \\
\texttt{But wait, there are 7 columns. So perhaps the separators are included in the string.} \\
\\ 
\texttt{Actually, the grid is drawn with the box characters, and the cells are the spaces. The string for row 1 has 7 characters? Let's count:} \\
\\ 
\texttt{"║◌MS│◌│◌║" {-} this has more than 7 characters. Probably it's represented with the box characters.} \\
\\ 
\texttt{To avoid confusion, I should rely on the coordinates. The grid has x from 1 to 7, y from 1 to 7.} \\
\\ 
\texttt{The proposed positions are (1,4) to (7,4). In my grid, row 4 is: `║◌◌◌│◌◌◌║ 4`} \\
\texttt{This means:} \\
\texttt{{-} Columns 1,2,3: all '◌' (so empty)} \\
\texttt{{-} Then a vertical line (not a cell)} \\
\texttt{{-} Columns 4,5,6: all '◌'} \\
\texttt{73? Wait, there should be 7 columns.} \\
\\ 
\texttt{Actually, row 4 has "◌◌◌│◌◌◌", which is 6 characters, but including the separators. The actual cells are at:} \\
\texttt{col1: first '◌'} \\
\texttt{col2: second '◌'} \\
\texttt{col3: third '◌'} \\
\texttt{col4: the first '◌' after the line? But the line is at the boundary.} \\
\\ 
\texttt{This is messy. To simplify, I should assume that all cells in row 4 are empty, as they are all shown as '◌'. So (1,4) to (7,4) should be empty.} \\
\\ 
\texttt{But note: the grid has fixed obstacles: the lines and corners. We cannot place objects on them. But the cells (1,4) to (7,4) are valid empty cells.} \\
\\ 
\texttt{Therefore, the proposed goal is acceptable.} \\
\\ 
\texttt{I should respond to the other player to agree.} \\
\\ 
\texttt{So my first command should be a SAY to acknowledge.} \\
\\ 
\texttt{I'll send:} \\
\texttt{SAY: I agree to your proposed goal: A(1,4), D(2,4), M(3,4), O(4,4), P(5,4), S(6,4), T(7,4). I will start moving objects to these positions.} \\
\\ 
\texttt{This way, we are synchronized.} \\
\\ 
\texttt{Therefore, my first command is:} \\
\texttt{```json} \\
\texttt{\{} \\
\texttt{  "command": "SAY: I agree to your proposed goal: A(1,4), D(2,4), M(3,4), O(4,4), P(5,4), S(6,4), T(7,4). I will start moving objects to these positions."} \\
\texttt{\}} \\
\texttt{```} \\
            }
        }
    }
     \\ \\

    \theutterance \stepcounter{utterance}  
    & & & \multicolumn{2}{p{0.3\linewidth}}{
        \cellcolor[rgb]{0.9,0.9,0.9}{
            \makecell[{{p{\linewidth}}}]{
                \texttt{\tiny{[GM$|$GM]}}
                \texttt{Let me analyze the situation. The other player has proposed a common goal configuration: all objects in row 4 from left to right: A at (1,4), D at (2,4), M at (3,4), O at (4,4), P at (5,4), S at (6,4), T at (7,4).} \\
\\ 
\texttt{However, looking at my grid, I can see that row 4 currently has three '◌' symbols and no objects. Specifically, row 4 is: `║◌◌◌│◌◌◌║ 4`. This means there are no objects in row 4 right now. Also, the grid has boundaries (like lines and corners) that might interfere with placements.} \\
\\ 
\texttt{But importantly, the proposed positions (1,4) to (7,4) are all within the grid and should be empty in my board (since row 4 is entirely empty of objects). So it's feasible.} \\
\\ 
\texttt{However, I should verify if all these target cells are indeed empty. Let's check:} \\
\\ 
\texttt{{-} (1,4): column 1, row 4 {-} this is the first cell in row 4, which is '◌' (empty).} \\
\texttt{{-} (2,4): also '◌'} \\
\texttt{{-} (3,4): also '◌'} \\
\texttt{{-} (4,4): this is the middle, but note there is a vertical line at column 4? Actually, looking at the grid, the columns are separated by lines. For example, in row 1: `║◌MS│◌│◌║` {-} the vertical lines are at column boundaries.} \\
\\ 
\texttt{Wait, let's understand the grid layout:} \\
\\ 
\texttt{The grid has 7 columns and 7 rows. The vertical lines are drawn between columns? Actually, the grid is drawn with box{-}drawing characters.} \\
\\ 
\texttt{Looking at row 1: `║◌MS│◌│◌║ 1`} \\
\texttt{This means:} \\
\texttt{{-} Column 1: ◌} \\
\texttt{{-} Column 2: M} \\
\texttt{{-} Column 3: S} \\
\texttt{{-} Column 4: │ (this is a line, not a cell)} \\
\texttt{Wait no, the lines are separators. The actual cells are between the lines.} \\
\\ 
\texttt{Actually, the grid is represented with the cells being the spaces between the lines. The numbers on top (1234567) indicate the column numbers.} \\
\\ 
\texttt{So for row 1:} \\
\texttt{{-} The left boundary is at x=0, but column 1 is at x=1.} \\
\texttt{{-} The string "◌MS│◌│◌" might be misleading.} \\
\\ 
\texttt{Let's parse the grid properly:} \\
\\ 
\texttt{The grid is 7x7. The top line "╔═══╤═╤═╗" indicates the top border, with openings.} \\
\\ 
\texttt{Better to think: each cell is one character wide? But some cells have multiple characters? Actually, no, each cell is one symbol.} \\
\\ 
\texttt{Actually, the grid is drawn with each cell being one character. The vertical lines are drawn between cells.} \\
\\ 
\texttt{For row 1: `║◌MS│◌│◌║` {-} this is confusing. Probably it's:} \\
\texttt{{-} First cell (col1): '◌'} \\
\texttt{{-} Second cell (col2): 'M'} \\
\texttt{{-} Third cell (col3): 'S'} \\
\texttt{{-} Then a vertical line '│' (which is a separator, not a cell)} \\
\texttt{{-} Then col4: '◌'} \\
\texttt{{-} Then a vertical line '│'} \\
\texttt{{-} Then col5: '◌'} \\
\texttt{But wait, there are 7 columns. So perhaps the separators are included in the string.} \\
\\ 
\texttt{Actually, the grid is drawn with the box characters, and the cells are the spaces. The string for row 1 has 7 characters? Let's count:} \\
\\ 
\texttt{"║◌MS│◌│◌║" {-} this has more than 7 characters. Probably it's represented with the box characters.} \\
\\ 
\texttt{To avoid confusion, I should rely on the coordinates. The grid has x from 1 to 7, y from 1 to 7.} \\
\\ 
\texttt{The proposed positions are (1,4) to (7,4). In my grid, row 4 is: `║◌◌◌│◌◌◌║ 4`} \\
\texttt{This means:} \\
\texttt{{-} Columns 1,2,3: all '◌' (so empty)} \\
\texttt{{-} Then a vertical line (not a cell)} \\
\texttt{{-} Columns 4,5,6: all '◌'} \\
\texttt{73? Wait, there should be 7 columns.} \\
\\ 
\texttt{Actually, row 4 has "◌◌◌│◌◌◌", which is 6 characters, but including the separators. The actual cells are at:} \\
\texttt{col1: first '◌'} \\
\texttt{col2: second '◌'} \\
\texttt{col3: third '◌'} \\
\texttt{col4: the first '◌' after the line? But the line is at the boundary.} \\
\\ 
\texttt{This is messy. To simplify, I should assume that all cells in row 4 are empty, as they are all shown as '◌'. So (1,4) to (7,4) should be empty.} \\
\\ 
\texttt{But note: the grid has fixed obstacles: the lines and corners. We cannot place objects on them. But the cells (1,4) to (7,4) are valid empty cells.} \\
\\ 
\texttt{Therefore, the proposed goal is acceptable.} \\
\\ 
\texttt{I should respond to the other player to agree.} \\
\\ 
\texttt{So my first command should be a SAY to acknowledge.} \\
\\ 
\texttt{I'll send:} \\
\texttt{SAY: I agree to your proposed goal: A(1,4), D(2,4), M(3,4), O(4,4), P(5,4), S(6,4), T(7,4). I will start moving objects to these positions.} \\
\\ 
\texttt{This way, we are synchronized.} \\
\\ 
\texttt{Therefore, my first command is:} \\
\texttt{```json} \\
\texttt{\{} \\
\texttt{  "command": "SAY: I agree to your proposed goal: A(1,4), D(2,4), M(3,4), O(4,4), P(5,4), S(6,4), T(7,4). I will start moving objects to these positions."} \\
\texttt{\}} \\
\texttt{```} \\
            }
        }
    }
    & & \\ \\

    \theutterance \stepcounter{utterance}  
    & & & \multicolumn{2}{p{0.3\linewidth}}{
        \cellcolor[rgb]{0.9,0.9,0.9}{
            \makecell[{{p{\linewidth}}}]{
                \texttt{\tiny{[GM$|$GM]}}
                \texttt{Invalid format: head and tail are not empty} \\
\texttt{head: 'Let me analyze the situation. The other player has proposed a common goal configuration: all objects in row 4 from left to right: A at (1,4), D at (2,4), M at (3,4), O at (4,4), P at (5,4), S at (6,4), T at (7,4).} \\
\\ 
\texttt{However, looking at my grid, I can see that row 4 currently has three '◌' symbols and no objects. Specifically, row 4 is: ║◌◌◌│◌◌◌║ 4. This means there are no objects in row 4 right now. Also, the grid has boundaries (like lines and corners) that might interfere with placements.} \\
\\ 
\texttt{But importantly, the proposed positions (1,4) to (7,4) are all within the grid and should be empty in my board (since row 4 is entirely empty of objects). So it's feasible.} \\
\\ 
\texttt{However, I should verify if all these target cells are indeed empty. Let's check:} \\
\\ 
\texttt{{-} (1,4): column 1, row 4 {-} this is the first cell in row 4, which is '◌' (empty).} \\
\texttt{{-} (2,4): also '◌'} \\
\texttt{{-} (3,4): also '◌'} \\
\texttt{{-} (4,4): this is the middle, but note there is a vertical line at column 4? Actually, looking at the grid, the columns are separated by lines. For example, in row 1: ║◌MS│◌│◌║ {-} the vertical lines are at column boundaries.} \\
\\ 
\texttt{Wait, let's understand the grid layout:} \\
\\ 
\texttt{The grid has 7 columns and 7 rows. The vertical lines are drawn between columns? Actually, the grid is drawn with box{-}drawing characters.} \\
\\ 
\texttt{Looking at row 1: ║◌MS│◌│◌║ 1} \\
\texttt{This means:} \\
\texttt{{-} Column 1: ◌} \\
\texttt{{-} Column 2: M} \\
\texttt{{-} Column 3: S} \\
\texttt{{-} Column 4: │ (this is a line, not a cell)} \\
\texttt{Wait no, the lines are separators. The actual cells are between the lines.} \\
\\ 
\texttt{Actually, the grid is represented with the cells being the spaces between the lines. The numbers on top (1234567) indicate the column numbers.} \\
\\ 
\texttt{So for row 1:} \\
\texttt{{-} The left boundary is at x=0, but column 1 is at x=1.} \\
\texttt{{-} The string "◌MS│◌│◌" might be misleading.} \\
\\ 
\texttt{Let's parse the grid properly:} \\
\\ 
\texttt{The grid is 7x7. The top line "╔═══╤═╤═╗" indicates the top border, with openings.} \\
\\ 
\texttt{Better to think: each cell is one character wide? But some cells have multiple characters? Actually, no, each cell is one symbol.} \\
\\ 
\texttt{Actually, the grid is drawn with each cell being one character. The vertical lines are drawn between cells.} \\
\\ 
\texttt{For row 1: ║◌MS│◌│◌║ {-} this is confusing. Probably it's:} \\
\texttt{{-} First cell (col1): '◌'} \\
\texttt{{-} Second cell (col2): 'M'} \\
\texttt{{-} Third cell (col3): 'S'} \\
\texttt{{-} Then a vertical line '│' (which is a separator, not a cell)} \\
\texttt{{-} Then col4: '◌'} \\
\texttt{{-} Then a vertical line '│'} \\
\texttt{{-} Then col5: '◌'} \\
\texttt{But wait, there are 7 columns. So perhaps the separators are included in the string.} \\
\\ 
\texttt{Actually, the grid is drawn with the box characters, and the cells are the spaces. The string for row 1 has 7 characters? Let's count:} \\
\\ 
\texttt{"║◌MS│◌│◌║" {-} this has more than 7 characters. Probably it's represented with the box characters.} \\
\\ 
\texttt{To avoid confusion, I should rely on the coordinates. The grid has x from 1 to 7, y from 1 to 7.} \\
\\ 
\texttt{The proposed positions are (1,4) to (7,4). In my grid, row 4 is: ║◌◌◌│◌◌◌║ 4} \\
\texttt{This means:} \\
\texttt{{-} Columns 1,2,3: all '◌' (so empty)} \\
\texttt{{-} Then a vertical line (not a cell)} \\
\texttt{{-} Columns 4,5,6: all '◌'} \\
\texttt{73? Wait, there should be 7 columns.} \\
\\ 
\texttt{Actually, row 4 has "◌◌◌│◌◌◌", which is 6 characters, but including the separators. The actual cells are at:} \\
\texttt{col1: first '◌'} \\
\texttt{col2: second '◌'} \\
\texttt{col3: third '◌'} \\
\texttt{col4: the first '◌' after the line? But the line is at the boundary.} \\
\\ 
\texttt{This is messy. To simplify, I should assume that all cells in row 4 are empty, as they are all shown as '◌'. So (1,4) to (7,4) should be empty.} \\
\\ 
\texttt{But note: the grid has fixed obstacles: the lines and corners. We cannot place objects on them. But the cells (1,4) to (7,4) are valid empty cells.} \\
\\ 
\texttt{Therefore, the proposed goal is acceptable.} \\
\\ 
\texttt{I should respond to the other player to agree.} \\
\\ 
\texttt{So my first command should be a SAY to acknowledge.} \\
\\ 
\texttt{I'll send:} \\
\texttt{SAY: I agree to your proposed goal: A(1,4), D(2,4), M(3,4), O(4,4), P(5,4), S(6,4), T(7,4). I will start moving objects to these positions.} \\
\\ 
\texttt{This way, we are synchronized.} \\
\\ 
\texttt{Therefore, my first command is:} \\
\texttt{json} \\
\texttt{\{} \\
\texttt{  "command": "'} \\
\texttt{tail: '} \\
\texttt{\}'} \\
            }
        }
    }
    & & \\ \\

    \theutterance \stepcounter{utterance}  
    & & \multicolumn{4}{p{0.6\linewidth}}{
        \cellcolor[rgb]{0.9,0.9,0.9}{
            \makecell[{{p{\linewidth}}}]{
                \texttt{\tiny{[GM$\rangle$P2]}}
                \texttt{Penalty: Your message must not contain anything before or after the command!} \\
\\ 
\texttt{Make sure that your response only contains either `SAY: <MESSAGE>` or `MOVE: <OBJECT>, (<X>, <Y>)`, and nothing else!} \\
\\ 
\texttt{You have collectively accumulated 1 of 16 penalties.} \\
\texttt{Please try again!} \\
            }
        }
    }
    & \\ \\

    \theutterance \stepcounter{utterance}  
    & & & \multicolumn{4}{p{0.6\linewidth}}{
        \cellcolor[rgb]{0.9,0.9,0.9}{
            \makecell[{{p{\linewidth}}}]{
                \texttt{\tiny{[GM$\langle$P2]}}
                \texttt{SAY: I agree to your proposed goal: A(1,4), D(2,4), M(3,4), O(4,4), P(5,4), S(6,4), T(7,4). I will start moving objects to these positions.} \\
            }
        }
    }
     \\ \\

    \theutterance \stepcounter{utterance}  
    & & & \multicolumn{2}{p{0.3\linewidth}}{
        \cellcolor[rgb]{0.9,0.9,0.9}{
            \makecell[{{p{\linewidth}}}]{
                \texttt{\tiny{[GM$|$GM]}}
                \texttt{SAY: I agree to your proposed goal: A(1,4), D(2,4), M(3,4), O(4,4), P(5,4), S(6,4), T(7,4). I will start moving objects to these positions.} \\
            }
        }
    }
    & & \\ \\

    \theutterance \stepcounter{utterance}  
    & & \multicolumn{4}{p{0.6\linewidth}}{
        \cellcolor[rgb]{0.9,0.9,0.9}{
            \makecell[{{p{\linewidth}}}]{
                \texttt{\tiny{[P1$\langle$GM]}}
                \texttt{Your message has been relayed to the other player.} \\
\\ 
\\ 
\texttt{You are currently playing round 1 of maximum 28.} \\
\texttt{You have collectively accumulated 1 of 16 penalties.} \\
\texttt{The other player sent this message:} \\
\\ 
\\ 
\texttt{"I agree to your proposed goal: A(1,4), D(2,4), M(3,4), O(4,4), P(5,4), S(6,4), T(7,4). I will start moving objects to these positions."} \\
\\ 
\\ 
\texttt{What is your next command?} \\
            }
        }
    }
    & \\ \\

    \theutterance \stepcounter{utterance}  
    & \multicolumn{4}{p{0.6\linewidth}}{
        \cellcolor[rgb]{0.9,0.9,0.9}{
            \makecell[{{p{\linewidth}}}]{
                \texttt{\tiny{[P1$\rangle$GM]}}
                \texttt{SAY: Great! I'll also start moving objects to the goal positions. Let me know if you encounter any issues. I'll begin with moving A to (1,4).} \\
            }
        }
    }
    & & \\ \\

    \theutterance \stepcounter{utterance}  
    & & & \multicolumn{2}{p{0.3\linewidth}}{
        \cellcolor[rgb]{0.9,0.9,0.9}{
            \makecell[{{p{\linewidth}}}]{
                \texttt{\tiny{[GM$|$GM]}}
                \texttt{SAY: Great! I'll also start moving objects to the goal positions. Let me know if you encounter any issues. I'll begin with moving A to (1,4).} \\
            }
        }
    }
    & & \\ \\

    \theutterance \stepcounter{utterance}  
    & & \multicolumn{4}{p{0.6\linewidth}}{
        \cellcolor[rgb]{0.9,0.9,0.9}{
            \makecell[{{p{\linewidth}}}]{
                \texttt{\tiny{[GM$\rangle$P2]}}
                \texttt{Your message has been relayed to the other player.} \\
\\ 
\\ 
\texttt{You are currently playing round 2 of maximum 28.} \\
\texttt{You have collectively accumulated 1 of 16 penalties.} \\
\texttt{The other player sent this message:} \\
\\ 
\\ 
\texttt{"Great! I'll also start moving objects to the goal positions. Let me know if you encounter any issues. I'll begin with moving A to (1,4)."} \\
\\ 
\\ 
\texttt{What is your next command?} \\
            }
        }
    }
    & \\ \\

    \theutterance \stepcounter{utterance}  
    & & & \multicolumn{4}{p{0.6\linewidth}}{
        \cellcolor[rgb]{0.9,0.9,0.9}{
            \makecell[{{p{\linewidth}}}]{
                \texttt{\tiny{[GM$\langle$P2]}}
                \texttt{Looking at my grid, I need to identify where the objects are currently located. From the grid:} \\
\\ 
\texttt{{-} Row 1: `║◌MS│◌│◌║` → This shows:} \\
\texttt{  {-} Column 1: ◌ (empty)} \\
\texttt{  {-} Column 2: M} \\
\texttt{  {-} Column 3: S} \\
\texttt{  {-} Column 4: │ (line, not a cell)} \\
\texttt{  {-} Column 5: ◌} \\
\texttt{  {-} Column 6: │ (line, not a cell)} \\
\texttt{  {-} Column 7: ◌} \\
\\ 
\texttt{But note: the grid has 7 columns. The vertical lines are separators between columns. Actually, the columns are:} \\
\\ 
\texttt{{-} Column 1: at position (1,1) is ◌} \\
\texttt{{-} Column 2: (2,1) is M} \\
\texttt{{-} Column 3: (3,1) is S} \\
\texttt{{-} Column 4: (4,1) is ◌ (after the first line? Wait, the line is drawn at the edge of column 3 and 4? This is confusing.} \\
\\ 
\texttt{To avoid confusion, I should refer to the coordinates as per the numbering. The top row has x from 1 to 7.} \\
\\ 
\texttt{In row 1:} \\
\texttt{{-} (1,1): ◌} \\
\texttt{{-} (2,1): M} \\
\texttt{{-} (3,1): S} \\
\texttt{{-} (4,1): this is where the line is? But the line is not a cell. Actually, the cell (4,1) is the one after the first vertical line. In the string "◌MS│◌│◌", the "│" is a separator. So:} \\
\texttt{  {-} After the left border "║", we have:} \\
\texttt{    {-} Cell (1,1): "◌"} \\
\texttt{    {-} Cell (2,1): "M"} \\
\texttt{    {-} Cell (3,1): "S"} \\
\texttt{    {-} Then a separator "│"} \\
\texttt{    {-} Cell (4,1): "◌"} \\
\texttt{    {-} Then a separator "│"} \\
\texttt{    {-} Cell (5,1): "◌"} \\
\texttt{    {-} Then the right border "║"} \\
\\ 
\texttt{But there are 7 columns. What about columns 6 and 7? The string might be abbreviated. Actually, the grid might have cells only at certain positions.} \\
\\ 
\texttt{Look at the entire grid:} \\
\\ 
\texttt{```} \\
\texttt{ 1234567} \\
\texttt{╔═══╤═╤═╗} \\
\texttt{║◌MS│◌│◌║ 1} \\
\texttt{║◌O◌│◌├─╢ 2} \\
\texttt{║◌◌◌├─┘◌║ 3} \\
\texttt{║◌◌◌│◌◌◌║ 4} \\
\texttt{╟─┬─┼───╢ 5} \\
\texttt{║◌│◌│◌◌D║ 6} \\
\texttt{║T│◌│◌PA║ 7} \\
\texttt{╚═╧═╧═══╝} \\
\texttt{```} \\
\\ 
\texttt{The top row has x{-}coordinates 1 to 7. The grid is drawn with box characters, so the cells are:} \\
\\ 
\texttt{{-} Row 1: cells at x=1,2,3,4,5,6,7? But the string "◌MS│◌│◌" has only 5 visible cells? This is not matching.} \\
\\ 
\texttt{Perhaps the "M" and "S" are in one cell? But no, each cell should contain one symbol.} \\
\\ 
\texttt{Another way: the grid has 7 columns, and the drawing uses multiple characters per cell? But that would be inefficient.} \\
\\ 
\texttt{Let's count the characters in the top row string: "║◌MS│◌│◌║" — let's break it:} \\
\\ 
\texttt{{-} "║" (1 char) {-} left border} \\
\texttt{{-} "◌" (1 char) {-} cell (1,1)} \\
\texttt{{-} "M" (1 char) {-} cell (2,1)} \\
\texttt{{-} "S" (1 char) {-} cell (3,1)} \\
\texttt{{-} "│" (1 char) {-} separator} \\
\texttt{{-} "◌" (1 char) {-} cell (4,1)} \\
\texttt{{-} "│" (1 char) {-} separator} \\
\texttt{{-} "◌" (1 char) {-} cell (5,1)} \\
\texttt{{-} "║" (1 char) {-} right border} \\
\\ 
\texttt{So only columns 1 to 5 are shown in this row? But there should be 7 columns. The grid might have variable width.} \\
\\ 
\texttt{Look at the top border: "╔═══╤═╤═╗" — this suggests that there are groups of columns. The "═══" might represent three columns, then "╤" is a junction, then "═" for one column, then "╤", then "═" for one column, but that doesn't add to 7.} \\
\\ 
\texttt{This is ambiguous. To resolve, I should assume that the grid has 7 columns, and the drawing is schematic. The objects are placed in the cells marked with letters.} \\
\\ 
\texttt{From the grid, I can see the following objects:} \\
\\ 
\texttt{{-} In row 1: M at (2,1), S at (3,1)} \\
\texttt{{-} In row 2: "║◌O◌│◌├─╢" — so O at (2,2)? And (1,2) is ◌, (3,2) is ◌, then separator, then (4,2) is ◌, then "├─" which might be part of the drawing.} \\
\texttt{{-} In row 3: "║◌◌◌├─┘◌║" — probably all empty except maybe at the end.} \\
\texttt{{-} In row 4: "║◌◌◌│◌◌◌║" — all empty.} \\
\texttt{{-} In row 6: "║◌│◌│◌◌D║" — D at (7,6) probably.} \\
\texttt{{-} In row 7: "║T│◌│◌PA║" — T at (1,7), P at (6,7), A at (7,7).} \\
\\ 
\texttt{Also, from the grid, there is 'O' in row 2, and 'T' in row 7, 'S' in row 1, 'M' in row 1, 'D' in row 6, 'P' and 'A' in row 7.} \\
\\ 
\texttt{So the objects are at:} \\
\texttt{{-} M: (2,1)} \\
\texttt{{-} S: (3,1)} \\
\texttt{{-} O: (2,2)} \\
\texttt{{-} D: (7,6)} \\
\texttt{{-} T: (1,7)} \\
\texttt{{-} P: (6,7)} \\
\texttt{{-} A: (7,7)} \\
\\ 
\texttt{Now, the goal is to move them to row 4:} \\
\texttt{{-} A to (1,4)} \\
\texttt{{-} D to (2,4)} \\
\texttt{{-} M to (3,4)} \\
\texttt{{-} O to (4,4)} \\
\texttt{{-} P to (5,4)} \\
\texttt{{-} S to (6,4)} \\
\texttt{{-} T to (7,4)} \\
\\ 
\texttt{I should start moving one object. Since the other player is moving A to (1,4), I should also move A to (1,4). But first, I need to check if (1,4) is empty.} \\
\\ 
\texttt{From row 4: "║◌◌◌│◌◌◌║" — this means:} \\
\texttt{{-} (1,4): ◌ (empty)} \\
\texttt{{-} (2,4): ◌} \\
\texttt{{-} (3,4): ◌} \\
\texttt{{-} (4,4): ◌ (after the line)} \\
\texttt{{-} (5,4): ◌} \\
\texttt{{-} (6,4): ◌} \\
\texttt{{-} (7,4): ◌} \\
\\ 
\texttt{So (1,4) is empty.} \\
\\ 
\texttt{Now, where is A on my board? From above, A is at (7,7).} \\
\\ 
\texttt{So I should move A from (7,7) to (1,4).} \\
\\ 
\texttt{Therefore, my next command is to move A to (1,4).} \\
\\ 
\texttt{```json} \\
\texttt{\{} \\
\texttt{  "command": "MOVE: A, (1,4)"} \\
\texttt{\}} \\
\texttt{```} \\
            }
        }
    }
     \\ \\

    \theutterance \stepcounter{utterance}  
    & & & \multicolumn{2}{p{0.3\linewidth}}{
        \cellcolor[rgb]{0.9,0.9,0.9}{
            \makecell[{{p{\linewidth}}}]{
                \texttt{\tiny{[GM$|$GM]}}
                \texttt{Looking at my grid, I need to identify where the objects are currently located. From the grid:} \\
\\ 
\texttt{{-} Row 1: `║◌MS│◌│◌║` → This shows:} \\
\texttt{  {-} Column 1: ◌ (empty)} \\
\texttt{  {-} Column 2: M} \\
\texttt{  {-} Column 3: S} \\
\texttt{  {-} Column 4: │ (line, not a cell)} \\
\texttt{  {-} Column 5: ◌} \\
\texttt{  {-} Column 6: │ (line, not a cell)} \\
\texttt{  {-} Column 7: ◌} \\
\\ 
\texttt{But note: the grid has 7 columns. The vertical lines are separators between columns. Actually, the columns are:} \\
\\ 
\texttt{{-} Column 1: at position (1,1) is ◌} \\
\texttt{{-} Column 2: (2,1) is M} \\
\texttt{{-} Column 3: (3,1) is S} \\
\texttt{{-} Column 4: (4,1) is ◌ (after the first line? Wait, the line is drawn at the edge of column 3 and 4? This is confusing.} \\
\\ 
\texttt{To avoid confusion, I should refer to the coordinates as per the numbering. The top row has x from 1 to 7.} \\
\\ 
\texttt{In row 1:} \\
\texttt{{-} (1,1): ◌} \\
\texttt{{-} (2,1): M} \\
\texttt{{-} (3,1): S} \\
\texttt{{-} (4,1): this is where the line is? But the line is not a cell. Actually, the cell (4,1) is the one after the first vertical line. In the string "◌MS│◌│◌", the "│" is a separator. So:} \\
\texttt{  {-} After the left border "║", we have:} \\
\texttt{    {-} Cell (1,1): "◌"} \\
\texttt{    {-} Cell (2,1): "M"} \\
\texttt{    {-} Cell (3,1): "S"} \\
\texttt{    {-} Then a separator "│"} \\
\texttt{    {-} Cell (4,1): "◌"} \\
\texttt{    {-} Then a separator "│"} \\
\texttt{    {-} Cell (5,1): "◌"} \\
\texttt{    {-} Then the right border "║"} \\
\\ 
\texttt{But there are 7 columns. What about columns 6 and 7? The string might be abbreviated. Actually, the grid might have cells only at certain positions.} \\
\\ 
\texttt{Look at the entire grid:} \\
\\ 
\texttt{```} \\
\texttt{ 1234567} \\
\texttt{╔═══╤═╤═╗} \\
\texttt{║◌MS│◌│◌║ 1} \\
\texttt{║◌O◌│◌├─╢ 2} \\
\texttt{║◌◌◌├─┘◌║ 3} \\
\texttt{║◌◌◌│◌◌◌║ 4} \\
\texttt{╟─┬─┼───╢ 5} \\
\texttt{║◌│◌│◌◌D║ 6} \\
\texttt{║T│◌│◌PA║ 7} \\
\texttt{╚═╧═╧═══╝} \\
\texttt{```} \\
\\ 
\texttt{The top row has x{-}coordinates 1 to 7. The grid is drawn with box characters, so the cells are:} \\
\\ 
\texttt{{-} Row 1: cells at x=1,2,3,4,5,6,7? But the string "◌MS│◌│◌" has only 5 visible cells? This is not matching.} \\
\\ 
\texttt{Perhaps the "M" and "S" are in one cell? But no, each cell should contain one symbol.} \\
\\ 
\texttt{Another way: the grid has 7 columns, and the drawing uses multiple characters per cell? But that would be inefficient.} \\
\\ 
\texttt{Let's count the characters in the top row string: "║◌MS│◌│◌║" — let's break it:} \\
\\ 
\texttt{{-} "║" (1 char) {-} left border} \\
\texttt{{-} "◌" (1 char) {-} cell (1,1)} \\
\texttt{{-} "M" (1 char) {-} cell (2,1)} \\
\texttt{{-} "S" (1 char) {-} cell (3,1)} \\
\texttt{{-} "│" (1 char) {-} separator} \\
\texttt{{-} "◌" (1 char) {-} cell (4,1)} \\
\texttt{{-} "│" (1 char) {-} separator} \\
\texttt{{-} "◌" (1 char) {-} cell (5,1)} \\
\texttt{{-} "║" (1 char) {-} right border} \\
\\ 
\texttt{So only columns 1 to 5 are shown in this row? But there should be 7 columns. The grid might have variable width.} \\
\\ 
\texttt{Look at the top border: "╔═══╤═╤═╗" — this suggests that there are groups of columns. The "═══" might represent three columns, then "╤" is a junction, then "═" for one column, then "╤", then "═" for one column, but that doesn't add to 7.} \\
\\ 
\texttt{This is ambiguous. To resolve, I should assume that the grid has 7 columns, and the drawing is schematic. The objects are placed in the cells marked with letters.} \\
\\ 
\texttt{From the grid, I can see the following objects:} \\
\\ 
\texttt{{-} In row 1: M at (2,1), S at (3,1)} \\
\texttt{{-} In row 2: "║◌O◌│◌├─╢" — so O at (2,2)? And (1,2) is ◌, (3,2) is ◌, then separator, then (4,2) is ◌, then "├─" which might be part of the drawing.} \\
\texttt{{-} In row 3: "║◌◌◌├─┘◌║" — probably all empty except maybe at the end.} \\
\texttt{{-} In row 4: "║◌◌◌│◌◌◌║" — all empty.} \\
\texttt{{-} In row 6: "║◌│◌│◌◌D║" — D at (7,6) probably.} \\
\texttt{{-} In row 7: "║T│◌│◌PA║" — T at (1,7), P at (6,7), A at (7,7).} \\
\\ 
\texttt{Also, from the grid, there is 'O' in row 2, and 'T' in row 7, 'S' in row 1, 'M' in row 1, 'D' in row 6, 'P' and 'A' in row 7.} \\
\\ 
\texttt{So the objects are at:} \\
\texttt{{-} M: (2,1)} \\
\texttt{{-} S: (3,1)} \\
\texttt{{-} O: (2,2)} \\
\texttt{{-} D: (7,6)} \\
\texttt{{-} T: (1,7)} \\
\texttt{{-} P: (6,7)} \\
\texttt{{-} A: (7,7)} \\
\\ 
\texttt{Now, the goal is to move them to row 4:} \\
\texttt{{-} A to (1,4)} \\
\texttt{{-} D to (2,4)} \\
\texttt{{-} M to (3,4)} \\
\texttt{{-} O to (4,4)} \\
\texttt{{-} P to (5,4)} \\
\texttt{{-} S to (6,4)} \\
\texttt{{-} T to (7,4)} \\
\\ 
\texttt{I should start moving one object. Since the other player is moving A to (1,4), I should also move A to (1,4). But first, I need to check if (1,4) is empty.} \\
\\ 
\texttt{From row 4: "║◌◌◌│◌◌◌║" — this means:} \\
\texttt{{-} (1,4): ◌ (empty)} \\
\texttt{{-} (2,4): ◌} \\
\texttt{{-} (3,4): ◌} \\
\texttt{{-} (4,4): ◌ (after the line)} \\
\texttt{{-} (5,4): ◌} \\
\texttt{{-} (6,4): ◌} \\
\texttt{{-} (7,4): ◌} \\
\\ 
\texttt{So (1,4) is empty.} \\
\\ 
\texttt{Now, where is A on my board? From above, A is at (7,7).} \\
\\ 
\texttt{So I should move A from (7,7) to (1,4).} \\
\\ 
\texttt{Therefore, my next command is to move A to (1,4).} \\
\\ 
\texttt{```json} \\
\texttt{\{} \\
\texttt{  "command": "MOVE: A, (1,4)"} \\
\texttt{\}} \\
\texttt{```} \\
            }
        }
    }
    & & \\ \\

    \theutterance \stepcounter{utterance}  
    & & & \multicolumn{2}{p{0.3\linewidth}}{
        \cellcolor[rgb]{0.9,0.9,0.9}{
            \makecell[{{p{\linewidth}}}]{
                \texttt{\tiny{[GM$|$GM]}}
                \texttt{Invalid format: head and tail are not empty} \\
\texttt{head: 'Looking at my grid, I need to identify where the objects are currently located. From the grid:} \\
\\ 
\texttt{{-} Row 1: ║◌MS│◌│◌║ → This shows:} \\
\texttt{  {-} Column 1: ◌ (empty)} \\
\texttt{  {-} Column 2: M} \\
\texttt{  {-} Column 3: S} \\
\texttt{  {-} Column 4: │ (line, not a cell)} \\
\texttt{  {-} Column 5: ◌} \\
\texttt{  {-} Column 6: │ (line, not a cell)} \\
\texttt{  {-} Column 7: ◌} \\
\\ 
\texttt{But note: the grid has 7 columns. The vertical lines are separators between columns. Actually, the columns are:} \\
\\ 
\texttt{{-} Column 1: at position (1,1) is ◌} \\
\texttt{{-} Column 2: (2,1) is M} \\
\texttt{{-} Column 3: (3,1) is S} \\
\texttt{{-} Column 4: (4,1) is ◌ (after the first line? Wait, the line is drawn at the edge of column 3 and 4? This is confusing.} \\
\\ 
\texttt{To avoid confusion, I should refer to the coordinates as per the numbering. The top row has x from 1 to 7.} \\
\\ 
\texttt{In row 1:} \\
\texttt{{-} (1,1): ◌} \\
\texttt{{-} (2,1): M} \\
\texttt{{-} (3,1): S} \\
\texttt{{-} (4,1): this is where the line is? But the line is not a cell. Actually, the cell (4,1) is the one after the first vertical line. In the string "◌MS│◌│◌", the "│" is a separator. So:} \\
\texttt{  {-} After the left border "║", we have:} \\
\texttt{    {-} Cell (1,1): "◌"} \\
\texttt{    {-} Cell (2,1): "M"} \\
\texttt{    {-} Cell (3,1): "S"} \\
\texttt{    {-} Then a separator "│"} \\
\texttt{    {-} Cell (4,1): "◌"} \\
\texttt{    {-} Then a separator "│"} \\
\texttt{    {-} Cell (5,1): "◌"} \\
\texttt{    {-} Then the right border "║"} \\
\\ 
\texttt{But there are 7 columns. What about columns 6 and 7? The string might be abbreviated. Actually, the grid might have cells only at certain positions.} \\
\\ 
\texttt{Look at the entire grid:} \\
\\ 
\\ 
\texttt{ 1234567} \\
\texttt{╔═══╤═╤═╗} \\
\texttt{║◌MS│◌│◌║ 1} \\
\texttt{║◌O◌│◌├─╢ 2} \\
\texttt{║◌◌◌├─┘◌║ 3} \\
\texttt{║◌◌◌│◌◌◌║ 4} \\
\texttt{╟─┬─┼───╢ 5} \\
\texttt{║◌│◌│◌◌D║ 6} \\
\texttt{║T│◌│◌PA║ 7} \\
\texttt{╚═╧═╧═══╝} \\
\\ 
\\ 
\texttt{The top row has x{-}coordinates 1 to 7. The grid is drawn with box characters, so the cells are:} \\
\\ 
\texttt{{-} Row 1: cells at x=1,2,3,4,5,6,7? But the string "◌MS│◌│◌" has only 5 visible cells? This is not matching.} \\
\\ 
\texttt{Perhaps the "M" and "S" are in one cell? But no, each cell should contain one symbol.} \\
\\ 
\texttt{Another way: the grid has 7 columns, and the drawing uses multiple characters per cell? But that would be inefficient.} \\
\\ 
\texttt{Let's count the characters in the top row string: "║◌MS│◌│◌║" — let's break it:} \\
\\ 
\texttt{{-} "║" (1 char) {-} left border} \\
\texttt{{-} "◌" (1 char) {-} cell (1,1)} \\
\texttt{{-} "M" (1 char) {-} cell (2,1)} \\
\texttt{{-} "S" (1 char) {-} cell (3,1)} \\
\texttt{{-} "│" (1 char) {-} separator} \\
\texttt{{-} "◌" (1 char) {-} cell (4,1)} \\
\texttt{{-} "│" (1 char) {-} separator} \\
\texttt{{-} "◌" (1 char) {-} cell (5,1)} \\
\texttt{{-} "║" (1 char) {-} right border} \\
\\ 
\texttt{So only columns 1 to 5 are shown in this row? But there should be 7 columns. The grid might have variable width.} \\
\\ 
\texttt{Look at the top border: "╔═══╤═╤═╗" — this suggests that there are groups of columns. The "═══" might represent three columns, then "╤" is a junction, then "═" for one column, then "╤", then "═" for one column, but that doesn't add to 7.} \\
\\ 
\texttt{This is ambiguous. To resolve, I should assume that the grid has 7 columns, and the drawing is schematic. The objects are placed in the cells marked with letters.} \\
\\ 
\texttt{From the grid, I can see the following objects:} \\
\\ 
\texttt{{-} In row 1: M at (2,1), S at (3,1)} \\
\texttt{{-} In row 2: "║◌O◌│◌├─╢" — so O at (2,2)? And (1,2) is ◌, (3,2) is ◌, then separator, then (4,2) is ◌, then "├─" which might be part of the drawing.} \\
\texttt{{-} In row 3: "║◌◌◌├─┘◌║" — probably all empty except maybe at the end.} \\
\texttt{{-} In row 4: "║◌◌◌│◌◌◌║" — all empty.} \\
\texttt{{-} In row 6: "║◌│◌│◌◌D║" — D at (7,6) probably.} \\
\texttt{{-} In row 7: "║T│◌│◌PA║" — T at (1,7), P at (6,7), A at (7,7).} \\
\\ 
\texttt{Also, from the grid, there is 'O' in row 2, and 'T' in row 7, 'S' in row 1, 'M' in row 1, 'D' in row 6, 'P' and 'A' in row 7.} \\
\\ 
\texttt{So the objects are at:} \\
\texttt{{-} M: (2,1)} \\
\texttt{{-} S: (3,1)} \\
\texttt{{-} O: (2,2)} \\
\texttt{{-} D: (7,6)} \\
\texttt{{-} T: (1,7)} \\
\texttt{{-} P: (6,7)} \\
\texttt{{-} A: (7,7)} \\
\\ 
\texttt{Now, the goal is to move them to row 4:} \\
\texttt{{-} A to (1,4)} \\
\texttt{{-} D to (2,4)} \\
\texttt{{-} M to (3,4)} \\
\texttt{{-} O to (4,4)} \\
\texttt{{-} P to (5,4)} \\
\texttt{{-} S to (6,4)} \\
\texttt{{-} T to (7,4)} \\
\\ 
\texttt{I should start moving one object. Since the other player is moving A to (1,4), I should also move A to (1,4). But first, I need to check if (1,4) is empty.} \\
\\ 
\texttt{From row 4: "║◌◌◌│◌◌◌║" — this means:} \\
\texttt{{-} (1,4): ◌ (empty)} \\
\texttt{{-} (2,4): ◌} \\
\texttt{{-} (3,4): ◌} \\
\texttt{{-} (4,4): ◌ (after the line)} \\
\texttt{{-} (5,4): ◌} \\
\texttt{{-} (6,4): ◌} \\
\texttt{{-} (7,4): ◌} \\
\\ 
\texttt{So (1,4) is empty.} \\
\\ 
\texttt{Now, where is A on my board? From above, A is at (7,7).} \\
\\ 
\texttt{So I should move A from (7,7) to (1,4).} \\
\\ 
\texttt{Therefore, my next command is to move A to (1,4).} \\
\\ 
\texttt{json} \\
\texttt{\{} \\
\texttt{  "command": "'} \\
\texttt{tail: '"} \\
\texttt{\}'} \\
            }
        }
    }
    & & \\ \\

    \theutterance \stepcounter{utterance}  
    & & \multicolumn{4}{p{0.6\linewidth}}{
        \cellcolor[rgb]{0.9,0.9,0.9}{
            \makecell[{{p{\linewidth}}}]{
                \texttt{\tiny{[GM$\rangle$P2]}}
                \texttt{Penalty: Your message must not contain anything before or after the command!} \\
\\ 
\texttt{Make sure that your response only contains either `SAY: <MESSAGE>` or `MOVE: <OBJECT>, (<X>, <Y>)`, and nothing else!} \\
\\ 
\texttt{You have collectively accumulated 2 of 16 penalties.} \\
\texttt{Please try again!} \\
            }
        }
    }
    & \\ \\

    \theutterance \stepcounter{utterance}  
    & & & \multicolumn{4}{p{0.6\linewidth}}{
        \cellcolor[rgb]{0.9,0.9,0.9}{
            \makecell[{{p{\linewidth}}}]{
                \texttt{\tiny{[GM$\langle$P2]}}
                \texttt{MOVE: A, (1,4)} \\
            }
        }
    }
     \\ \\

    \theutterance \stepcounter{utterance}  
    & & & \multicolumn{2}{p{0.3\linewidth}}{
        \cellcolor[rgb]{0.9,0.9,0.9}{
            \makecell[{{p{\linewidth}}}]{
                \texttt{\tiny{[GM$|$GM]}}
                \texttt{MOVE: A, (1,4)} \\
            }
        }
    }
    & & \\ \\

    \theutterance \stepcounter{utterance}  
    & & & \multicolumn{2}{p{0.3\linewidth}}{
        \cellcolor[rgb]{0.9,0.9,0.9}{
            \makecell[{{p{\linewidth}}}]{
                \texttt{\tiny{[GM$|$GM]}}
                \texttt{Moved 'A' to (1,4) successfully. Your updated grid looks like this:} \\
\texttt{```} \\
\texttt{ 1234567} \\
\texttt{╔═══╤═╤═╗} \\
\texttt{║◌MS│◌│◌║ 1} \\
\texttt{║◌O◌│◌├─╢ 2} \\
\texttt{║◌◌◌├─┘◌║ 3} \\
\texttt{║A◌◌│◌◌◌║ 4} \\
\texttt{╟─┬─┼───╢ 5} \\
\texttt{║◌│◌│◌◌D║ 6} \\
\texttt{║T│◌│◌P◌║ 7} \\
\texttt{╚═╧═╧═══╝} \\
\texttt{```} \\
            }
        }
    }
    & & \\ \\

    \theutterance \stepcounter{utterance}  
    & & \multicolumn{4}{p{0.6\linewidth}}{
        \cellcolor[rgb]{0.9,0.9,0.9}{
            \makecell[{{p{\linewidth}}}]{
                \texttt{\tiny{[P1$\langle$GM]}}
                \texttt{Your message has been relayed to the other player.} \\
\\ 
\\ 
\texttt{You are currently playing round 2 of maximum 28.} \\
\texttt{You have collectively accumulated 2 of 16 penalties.} \\
\texttt{The other player moved an object on their grid.} \\
\\ 
\\ 
\texttt{What is your next command?} \\
            }
        }
    }
    & \\ \\

    \theutterance \stepcounter{utterance}  
    & \multicolumn{4}{p{0.6\linewidth}}{
        \cellcolor[rgb]{0.9,0.9,0.9}{
            \makecell[{{p{\linewidth}}}]{
                \texttt{\tiny{[P1$\rangle$GM]}}
                \texttt{MOVE: A, (1,4)} \\
            }
        }
    }
    & & \\ \\

    \theutterance \stepcounter{utterance}  
    & & & \multicolumn{2}{p{0.3\linewidth}}{
        \cellcolor[rgb]{0.9,0.9,0.9}{
            \makecell[{{p{\linewidth}}}]{
                \texttt{\tiny{[GM$|$GM]}}
                \texttt{MOVE: A, (1,4)} \\
            }
        }
    }
    & & \\ \\

    \theutterance \stepcounter{utterance}  
    & & & \multicolumn{2}{p{0.3\linewidth}}{
        \cellcolor[rgb]{0.9,0.9,0.9}{
            \makecell[{{p{\linewidth}}}]{
                \texttt{\tiny{[GM$|$GM]}}
                \texttt{Moved 'A' to (1,4) successfully. Your updated grid looks like this:} \\
\texttt{```} \\
\texttt{ 1234567} \\
\texttt{╔═══╤═╤═╗} \\
\texttt{║◌◌◌│◌│◌║ 1} \\
\texttt{║◌◌M│◌├─╢ 2} \\
\texttt{║◌T◌├─┘◌║ 3} \\
\texttt{║A◌P│◌◌◌║ 4} \\
\texttt{╟─┬─┼───╢ 5} \\
\texttt{║◌│◌│S◌O║ 6} \\
\texttt{║◌│◌│◌D◌║ 7} \\
\texttt{╚═╧═╧═══╝} \\
\texttt{```} \\
            }
        }
    }
    & & \\ \\

    \theutterance \stepcounter{utterance}  
    & & \multicolumn{4}{p{0.6\linewidth}}{
        \cellcolor[rgb]{0.9,0.9,0.9}{
            \makecell[{{p{\linewidth}}}]{
                \texttt{\tiny{[GM$\rangle$P2]}}
                \texttt{Moved 'A' to (1,4) successfully. Your updated grid looks like this:} \\
\texttt{```} \\
\texttt{ 1234567} \\
\texttt{╔═══╤═╤═╗} \\
\texttt{║◌MS│◌│◌║ 1} \\
\texttt{║◌O◌│◌├─╢ 2} \\
\texttt{║◌◌◌├─┘◌║ 3} \\
\texttt{║A◌◌│◌◌◌║ 4} \\
\texttt{╟─┬─┼───╢ 5} \\
\texttt{║◌│◌│◌◌D║ 6} \\
\texttt{║T│◌│◌P◌║ 7} \\
\texttt{╚═╧═╧═══╝} \\
\texttt{```} \\
\texttt{You are currently playing round 3 of maximum 28.} \\
\texttt{You have collectively accumulated 2 of 16 penalties.} \\
\texttt{The other player moved an object on their grid.} \\
\\ 
\\ 
\texttt{What is your next command?} \\
            }
        }
    }
    & \\ \\

    \theutterance \stepcounter{utterance}  
    & & & \multicolumn{4}{p{0.6\linewidth}}{
        \cellcolor[rgb]{0.9,0.9,0.9}{
            \makecell[{{p{\linewidth}}}]{
                \texttt{\tiny{[GM$\langle$P2]}}
                \texttt{Looking at my updated grid, I have successfully moved A to (1,4). Now I should continue moving other objects to their goal positions. The next logical step is to move D to (2,4), but first I need to check where D is currently.} \\
\\ 
\texttt{From the grid:} \\
\texttt{{-} D is at (7,6) as seen in row 6: `║◌│◌│◌◌D║ 6`} \\
\\ 
\texttt{I should move D to (2,4). Check if (2,4) is empty: from row 4 `║A◌◌│◌◌◌║ 4`, column 2 is currently empty (◌).} \\
\\ 
\texttt{So, I will move D to (2,4).} \\
\\ 
\texttt{MOVE: D, (2,4)} \\
            }
        }
    }
     \\ \\

    \theutterance \stepcounter{utterance}  
    & & & \multicolumn{2}{p{0.3\linewidth}}{
        \cellcolor[rgb]{0.9,0.9,0.9}{
            \makecell[{{p{\linewidth}}}]{
                \texttt{\tiny{[GM$|$GM]}}
                \texttt{Looking at my updated grid, I have successfully moved A to (1,4). Now I should continue moving other objects to their goal positions. The next logical step is to move D to (2,4), but first I need to check where D is currently.} \\
\\ 
\texttt{From the grid:} \\
\texttt{{-} D is at (7,6) as seen in row 6: `║◌│◌│◌◌D║ 6`} \\
\\ 
\texttt{I should move D to (2,4). Check if (2,4) is empty: from row 4 `║A◌◌│◌◌◌║ 4`, column 2 is currently empty (◌).} \\
\\ 
\texttt{So, I will move D to (2,4).} \\
\\ 
\texttt{MOVE: D, (2,4)} \\
            }
        }
    }
    & & \\ \\

    \theutterance \stepcounter{utterance}  
    & & & \multicolumn{2}{p{0.3\linewidth}}{
        \cellcolor[rgb]{0.9,0.9,0.9}{
            \makecell[{{p{\linewidth}}}]{
                \texttt{\tiny{[GM$|$GM]}}
                \texttt{Invalid format: head is not empty: 'Looking at my updated grid, I have successfully moved A to (1,4). Now I should continue moving other objects to their goal positions. The next logical step is to move D to (2,4), but first I need to check where D is currently.} \\
\\ 
\texttt{From the grid:} \\
\texttt{{-} D is at (7,6) as seen in row 6: ║◌│◌│◌◌D║ 6} \\
\\ 
\texttt{I should move D to (2,4). Check if (2,4) is empty: from row 4 ║A◌◌│◌◌◌║ 4, column 2 is currently empty (◌).} \\
\\ 
\texttt{So, I will move D to (2,4).} \\
\\ 
\texttt{'} \\
            }
        }
    }
    & & \\ \\

    \theutterance \stepcounter{utterance}  
    & & \multicolumn{4}{p{0.6\linewidth}}{
        \cellcolor[rgb]{0.9,0.9,0.9}{
            \makecell[{{p{\linewidth}}}]{
                \texttt{\tiny{[GM$\rangle$P2]}}
                \texttt{Penalty: Your message must not contain anything before the command!} \\
\\ 
\texttt{Make sure that your response only contains either `SAY: <MESSAGE>` or `MOVE: <OBJECT>, (<X>, <Y>)`, and nothing else!} \\
\\ 
\texttt{You have collectively accumulated 3 of 16 penalties.} \\
\texttt{Please try again!} \\
            }
        }
    }
    & \\ \\

    \theutterance \stepcounter{utterance}  
    & & & \multicolumn{4}{p{0.6\linewidth}}{
        \cellcolor[rgb]{0.9,0.9,0.9}{
            \makecell[{{p{\linewidth}}}]{
                \texttt{\tiny{[GM$\langle$P2]}}
                \texttt{MOVE: D, (2,4)} \\
            }
        }
    }
     \\ \\

    \theutterance \stepcounter{utterance}  
    & & & \multicolumn{2}{p{0.3\linewidth}}{
        \cellcolor[rgb]{0.9,0.9,0.9}{
            \makecell[{{p{\linewidth}}}]{
                \texttt{\tiny{[GM$|$GM]}}
                \texttt{MOVE: D, (2,4)} \\
            }
        }
    }
    & & \\ \\

    \theutterance \stepcounter{utterance}  
    & & & \multicolumn{2}{p{0.3\linewidth}}{
        \cellcolor[rgb]{0.9,0.9,0.9}{
            \makecell[{{p{\linewidth}}}]{
                \texttt{\tiny{[GM$|$GM]}}
                \texttt{Moved 'D' to (2,4) successfully. Your updated grid looks like this:} \\
\texttt{```} \\
\texttt{ 1234567} \\
\texttt{╔═══╤═╤═╗} \\
\texttt{║◌MS│◌│◌║ 1} \\
\texttt{║◌O◌│◌├─╢ 2} \\
\texttt{║◌◌◌├─┘◌║ 3} \\
\texttt{║AD◌│◌◌◌║ 4} \\
\texttt{╟─┬─┼───╢ 5} \\
\texttt{║◌│◌│◌◌◌║ 6} \\
\texttt{║T│◌│◌P◌║ 7} \\
\texttt{╚═╧═╧═══╝} \\
\texttt{```} \\
            }
        }
    }
    & & \\ \\

    \theutterance \stepcounter{utterance}  
    & & \multicolumn{4}{p{0.6\linewidth}}{
        \cellcolor[rgb]{0.9,0.9,0.9}{
            \makecell[{{p{\linewidth}}}]{
                \texttt{\tiny{[P1$\langle$GM]}}
                \texttt{Moved 'A' to (1,4) successfully. Your updated grid looks like this:} \\
\texttt{```} \\
\texttt{ 1234567} \\
\texttt{╔═══╤═╤═╗} \\
\texttt{║◌◌◌│◌│◌║ 1} \\
\texttt{║◌◌M│◌├─╢ 2} \\
\texttt{║◌T◌├─┘◌║ 3} \\
\texttt{║A◌P│◌◌◌║ 4} \\
\texttt{╟─┬─┼───╢ 5} \\
\texttt{║◌│◌│S◌O║ 6} \\
\texttt{║◌│◌│◌D◌║ 7} \\
\texttt{╚═╧═╧═══╝} \\
\texttt{```} \\
\texttt{You are currently playing round 3 of maximum 28.} \\
\texttt{You have collectively accumulated 3 of 16 penalties.} \\
\texttt{The other player moved an object on their grid.} \\
\\ 
\\ 
\texttt{What is your next command?} \\
            }
        }
    }
    & \\ \\

    \theutterance \stepcounter{utterance}  
    & \multicolumn{4}{p{0.6\linewidth}}{
        \cellcolor[rgb]{0.9,0.9,0.9}{
            \makecell[{{p{\linewidth}}}]{
                \texttt{\tiny{[P1$\rangle$GM]}}
                \texttt{MOVE: D, (2,4)} \\
            }
        }
    }
    & & \\ \\

    \theutterance \stepcounter{utterance}  
    & & & \multicolumn{2}{p{0.3\linewidth}}{
        \cellcolor[rgb]{0.9,0.9,0.9}{
            \makecell[{{p{\linewidth}}}]{
                \texttt{\tiny{[GM$|$GM]}}
                \texttt{MOVE: D, (2,4)} \\
            }
        }
    }
    & & \\ \\

    \theutterance \stepcounter{utterance}  
    & & & \multicolumn{2}{p{0.3\linewidth}}{
        \cellcolor[rgb]{0.9,0.9,0.9}{
            \makecell[{{p{\linewidth}}}]{
                \texttt{\tiny{[GM$|$GM]}}
                \texttt{Moved 'D' to (2,4) successfully. Your updated grid looks like this:} \\
\texttt{```} \\
\texttt{ 1234567} \\
\texttt{╔═══╤═╤═╗} \\
\texttt{║◌◌◌│◌│◌║ 1} \\
\texttt{║◌◌M│◌├─╢ 2} \\
\texttt{║◌T◌├─┘◌║ 3} \\
\texttt{║ADP│◌◌◌║ 4} \\
\texttt{╟─┬─┼───╢ 5} \\
\texttt{║◌│◌│S◌O║ 6} \\
\texttt{║◌│◌│◌◌◌║ 7} \\
\texttt{╚═╧═╧═══╝} \\
\texttt{```} \\
            }
        }
    }
    & & \\ \\

    \theutterance \stepcounter{utterance}  
    & & \multicolumn{4}{p{0.6\linewidth}}{
        \cellcolor[rgb]{0.9,0.9,0.9}{
            \makecell[{{p{\linewidth}}}]{
                \texttt{\tiny{[GM$\rangle$P2]}}
                \texttt{Moved 'D' to (2,4) successfully. Your updated grid looks like this:} \\
\texttt{```} \\
\texttt{ 1234567} \\
\texttt{╔═══╤═╤═╗} \\
\texttt{║◌MS│◌│◌║ 1} \\
\texttt{║◌O◌│◌├─╢ 2} \\
\texttt{║◌◌◌├─┘◌║ 3} \\
\texttt{║AD◌│◌◌◌║ 4} \\
\texttt{╟─┬─┼───╢ 5} \\
\texttt{║◌│◌│◌◌◌║ 6} \\
\texttt{║T│◌│◌P◌║ 7} \\
\texttt{╚═╧═╧═══╝} \\
\texttt{```} \\
\texttt{You are currently playing round 4 of maximum 28.} \\
\texttt{You have collectively accumulated 3 of 16 penalties.} \\
\texttt{The other player moved an object on their grid.} \\
\\ 
\\ 
\texttt{What is your next command?} \\
            }
        }
    }
    & \\ \\

    \theutterance \stepcounter{utterance}  
    & & & \multicolumn{4}{p{0.6\linewidth}}{
        \cellcolor[rgb]{0.9,0.9,0.9}{
            \makecell[{{p{\linewidth}}}]{
                \texttt{\tiny{[GM$\langle$P2]}}
                \texttt{MOVE: M, (3,4)} \\
            }
        }
    }
     \\ \\

    \theutterance \stepcounter{utterance}  
    & & & \multicolumn{2}{p{0.3\linewidth}}{
        \cellcolor[rgb]{0.9,0.9,0.9}{
            \makecell[{{p{\linewidth}}}]{
                \texttt{\tiny{[GM$|$GM]}}
                \texttt{MOVE: M, (3,4)} \\
            }
        }
    }
    & & \\ \\

    \theutterance \stepcounter{utterance}  
    & & & \multicolumn{2}{p{0.3\linewidth}}{
        \cellcolor[rgb]{0.9,0.9,0.9}{
            \makecell[{{p{\linewidth}}}]{
                \texttt{\tiny{[GM$|$GM]}}
                \texttt{Moved 'M' to (3,4) successfully. Your updated grid looks like this:} \\
\texttt{```} \\
\texttt{ 1234567} \\
\texttt{╔═══╤═╤═╗} \\
\texttt{║◌◌S│◌│◌║ 1} \\
\texttt{║◌O◌│◌├─╢ 2} \\
\texttt{║◌◌◌├─┘◌║ 3} \\
\texttt{║ADM│◌◌◌║ 4} \\
\texttt{╟─┬─┼───╢ 5} \\
\texttt{║◌│◌│◌◌◌║ 6} \\
\texttt{║T│◌│◌P◌║ 7} \\
\texttt{╚═╧═╧═══╝} \\
\texttt{```} \\
            }
        }
    }
    & & \\ \\

    \theutterance \stepcounter{utterance}  
    & & \multicolumn{4}{p{0.6\linewidth}}{
        \cellcolor[rgb]{0.9,0.9,0.9}{
            \makecell[{{p{\linewidth}}}]{
                \texttt{\tiny{[P1$\langle$GM]}}
                \texttt{Moved 'D' to (2,4) successfully. Your updated grid looks like this:} \\
\texttt{```} \\
\texttt{ 1234567} \\
\texttt{╔═══╤═╤═╗} \\
\texttt{║◌◌◌│◌│◌║ 1} \\
\texttt{║◌◌M│◌├─╢ 2} \\
\texttt{║◌T◌├─┘◌║ 3} \\
\texttt{║ADP│◌◌◌║ 4} \\
\texttt{╟─┬─┼───╢ 5} \\
\texttt{║◌│◌│S◌O║ 6} \\
\texttt{║◌│◌│◌◌◌║ 7} \\
\texttt{╚═╧═╧═══╝} \\
\texttt{```} \\
\texttt{You are currently playing round 4 of maximum 28.} \\
\texttt{You have collectively accumulated 3 of 16 penalties.} \\
\texttt{The other player moved an object on their grid.} \\
\\ 
\\ 
\texttt{What is your next command?} \\
            }
        }
    }
    & \\ \\

    \theutterance \stepcounter{utterance}  
    & \multicolumn{4}{p{0.6\linewidth}}{
        \cellcolor[rgb]{0.9,0.9,0.9}{
            \makecell[{{p{\linewidth}}}]{
                \texttt{\tiny{[P1$\rangle$GM]}}
                \texttt{MOVE: M, (3,4)} \\
            }
        }
    }
    & & \\ \\

    \theutterance \stepcounter{utterance}  
    & & & \multicolumn{2}{p{0.3\linewidth}}{
        \cellcolor[rgb]{0.9,0.9,0.9}{
            \makecell[{{p{\linewidth}}}]{
                \texttt{\tiny{[GM$|$GM]}}
                \texttt{MOVE: M, (3,4)} \\
            }
        }
    }
    & & \\ \\

    \theutterance \stepcounter{utterance}  
    & & & \multicolumn{2}{p{0.3\linewidth}}{
        \cellcolor[rgb]{0.9,0.9,0.9}{
            \makecell[{{p{\linewidth}}}]{
                \texttt{\tiny{[GM$|$GM]}}
                \texttt{Penalty: (3,4) is not empty, but contains 'P'.} \\
\\ 
\texttt{You have collectively accumulated 4 of 16 penalties.} \\
\texttt{Please try again!} \\
            }
        }
    }
    & & \\ \\

    \theutterance \stepcounter{utterance}  
    & & \multicolumn{4}{p{0.6\linewidth}}{
        \cellcolor[rgb]{0.9,0.9,0.9}{
            \makecell[{{p{\linewidth}}}]{
                \texttt{\tiny{[P1$\langle$GM]}}
                \texttt{Penalty: (3,4) is not empty, but contains 'P'.} \\
\\ 
\texttt{You have collectively accumulated 4 of 16 penalties.} \\
\texttt{Please try again!} \\
            }
        }
    }
    & \\ \\

    \theutterance \stepcounter{utterance}  
    & \multicolumn{4}{p{0.6\linewidth}}{
        \cellcolor[rgb]{0.9,0.9,0.9}{
            \makecell[{{p{\linewidth}}}]{
                \texttt{\tiny{[P1$\rangle$GM]}}
                \texttt{SAY: I see P is at (3,4) now, so I need to move P first. I'll move P to (5,4) to free up (3,4) for M.} \\
            }
        }
    }
    & & \\ \\

    \theutterance \stepcounter{utterance}  
    & & & \multicolumn{2}{p{0.3\linewidth}}{
        \cellcolor[rgb]{0.9,0.9,0.9}{
            \makecell[{{p{\linewidth}}}]{
                \texttt{\tiny{[GM$|$GM]}}
                \texttt{SAY: I see P is at (3,4) now, so I need to move P first. I'll move P to (5,4) to free up (3,4) for M.} \\
            }
        }
    }
    & & \\ \\

    \theutterance \stepcounter{utterance}  
    & & \multicolumn{4}{p{0.6\linewidth}}{
        \cellcolor[rgb]{0.9,0.9,0.9}{
            \makecell[{{p{\linewidth}}}]{
                \texttt{\tiny{[GM$\rangle$P2]}}
                \texttt{Moved 'M' to (3,4) successfully. Your updated grid looks like this:} \\
\texttt{```} \\
\texttt{ 1234567} \\
\texttt{╔═══╤═╤═╗} \\
\texttt{║◌◌S│◌│◌║ 1} \\
\texttt{║◌O◌│◌├─╢ 2} \\
\texttt{║◌◌◌├─┘◌║ 3} \\
\texttt{║ADM│◌◌◌║ 4} \\
\texttt{╟─┬─┼───╢ 5} \\
\texttt{║◌│◌│◌◌◌║ 6} \\
\texttt{║T│◌│◌P◌║ 7} \\
\texttt{╚═╧═╧═══╝} \\
\texttt{```} \\
\texttt{You are currently playing round 5 of maximum 28.} \\
\texttt{You have collectively accumulated 4 of 16 penalties.} \\
\texttt{The other player sent this message:} \\
\\ 
\\ 
\texttt{"I see P is at (3,4) now, so I need to move P first. I'll move P to (5,4) to free up (3,4) for M."} \\
\\ 
\\ 
\texttt{What is your next command?} \\
            }
        }
    }
    & \\ \\

    \theutterance \stepcounter{utterance}  
    & & & \multicolumn{4}{p{0.6\linewidth}}{
        \cellcolor[rgb]{0.9,0.9,0.9}{
            \makecell[{{p{\linewidth}}}]{
                \texttt{\tiny{[GM$\langle$P2]}}
                \texttt{Looking at my grid, I have just moved M to (3,4), so it is correctly placed. The other player seems to have a different configuration — they mention that P is at (3,4) on their board, but on my board, M is at (3,4). This indicates that our boards are not identical, which is expected.} \\
\\ 
\texttt{The other player is moving their P from (3,4) to (5,4) to free up (3,4) for their M. On my board, I need to move O to (4,4). Currently, O is at (2,2) as seen in row 2: `║◌O◌│◌├─╢ 2`.} \\
\\ 
\texttt{I should move O to (4,4). Check row 4: `║ADM│◌◌◌║ 4` — the cell at (4,4) is after the vertical line, which is the first cell in the second group. So (4,4) is empty (◌).} \\
\\ 
\texttt{Therefore, I will move O to (4,4).} \\
\\ 
\texttt{MOVE: O, (4,4)} \\
            }
        }
    }
     \\ \\

    \theutterance \stepcounter{utterance}  
    & & & \multicolumn{2}{p{0.3\linewidth}}{
        \cellcolor[rgb]{0.9,0.9,0.9}{
            \makecell[{{p{\linewidth}}}]{
                \texttt{\tiny{[GM$|$GM]}}
                \texttt{Looking at my grid, I have just moved M to (3,4), so it is correctly placed. The other player seems to have a different configuration — they mention that P is at (3,4) on their board, but on my board, M is at (3,4). This indicates that our boards are not identical, which is expected.} \\
\\ 
\texttt{The other player is moving their P from (3,4) to (5,4) to free up (3,4) for their M. On my board, I need to move O to (4,4). Currently, O is at (2,2) as seen in row 2: `║◌O◌│◌├─╢ 2`.} \\
\\ 
\texttt{I should move O to (4,4). Check row 4: `║ADM│◌◌◌║ 4` — the cell at (4,4) is after the vertical line, which is the first cell in the second group. So (4,4) is empty (◌).} \\
\\ 
\texttt{Therefore, I will move O to (4,4).} \\
\\ 
\texttt{MOVE: O, (4,4)} \\
            }
        }
    }
    & & \\ \\

    \theutterance \stepcounter{utterance}  
    & & & \multicolumn{2}{p{0.3\linewidth}}{
        \cellcolor[rgb]{0.9,0.9,0.9}{
            \makecell[{{p{\linewidth}}}]{
                \texttt{\tiny{[GM$|$GM]}}
                \texttt{Invalid format: head is not empty: 'Looking at my grid, I have just moved M to (3,4), so it is correctly placed. The other player seems to have a different configuration — they mention that P is at (3,4) on their board, but on my board, M is at (3,4). This indicates that our boards are not identical, which is expected.} \\
\\ 
\texttt{The other player is moving their P from (3,4) to (5,4) to free up (3,4) for their M. On my board, I need to move O to (4,4). Currently, O is at (2,2) as seen in row 2: ║◌O◌│◌├─╢ 2.} \\
\\ 
\texttt{I should move O to (4,4). Check row 4: ║ADM│◌◌◌║ 4 — the cell at (4,4) is after the vertical line, which is the first cell in the second group. So (4,4) is empty (◌).} \\
\\ 
\texttt{Therefore, I will move O to (4,4).} \\
\\ 
\texttt{'} \\
            }
        }
    }
    & & \\ \\

    \theutterance \stepcounter{utterance}  
    & & \multicolumn{4}{p{0.6\linewidth}}{
        \cellcolor[rgb]{0.9,0.9,0.9}{
            \makecell[{{p{\linewidth}}}]{
                \texttt{\tiny{[GM$\rangle$P2]}}
                \texttt{Penalty: Your message must not contain anything before the command!} \\
\\ 
\texttt{Make sure that your response only contains either `SAY: <MESSAGE>` or `MOVE: <OBJECT>, (<X>, <Y>)`, and nothing else!} \\
\\ 
\texttt{You have collectively accumulated 5 of 16 penalties.} \\
\texttt{Please try again!} \\
            }
        }
    }
    & \\ \\

    \theutterance \stepcounter{utterance}  
    & & & \multicolumn{4}{p{0.6\linewidth}}{
        \cellcolor[rgb]{0.9,0.9,0.9}{
            \makecell[{{p{\linewidth}}}]{
                \texttt{\tiny{[GM$\langle$P2]}}
                \texttt{MOVE: O, (4,4)} \\
            }
        }
    }
     \\ \\

    \theutterance \stepcounter{utterance}  
    & & & \multicolumn{2}{p{0.3\linewidth}}{
        \cellcolor[rgb]{0.9,0.9,0.9}{
            \makecell[{{p{\linewidth}}}]{
                \texttt{\tiny{[GM$|$GM]}}
                \texttt{MOVE: O, (4,4)} \\
            }
        }
    }
    & & \\ \\

    \theutterance \stepcounter{utterance}  
    & & & \multicolumn{2}{p{0.3\linewidth}}{
        \cellcolor[rgb]{0.9,0.9,0.9}{
            \makecell[{{p{\linewidth}}}]{
                \texttt{\tiny{[GM$|$GM]}}
                \texttt{Penalty: (4,4) is not empty, but contains '│'.} \\
\\ 
\texttt{You have collectively accumulated 6 of 16 penalties.} \\
\texttt{Please try again!} \\
            }
        }
    }
    & & \\ \\

    \theutterance \stepcounter{utterance}  
    & & \multicolumn{4}{p{0.6\linewidth}}{
        \cellcolor[rgb]{0.9,0.9,0.9}{
            \makecell[{{p{\linewidth}}}]{
                \texttt{\tiny{[GM$\rangle$P2]}}
                \texttt{Penalty: (4,4) is not empty, but contains '│'.} \\
\\ 
\texttt{You have collectively accumulated 6 of 16 penalties.} \\
\texttt{Please try again!} \\
            }
        }
    }
    & \\ \\

    \theutterance \stepcounter{utterance}  
    & & & \multicolumn{4}{p{0.6\linewidth}}{
        \cellcolor[rgb]{0.9,0.9,0.9}{
            \makecell[{{p{\linewidth}}}]{
                \texttt{\tiny{[GM$\langle$P2]}}
                \texttt{I see the issue: (4,4) is a vertical line '│', not an empty cell. The grid uses box{-}drawing characters, so I must avoid those. Looking at row 4: `║ADM│◌◌◌║ 4`, the cells are:} \\
\\ 
\texttt{{-} (1,4): A} \\
\texttt{{-} (2,4): D} \\
\texttt{{-} (3,4): M} \\
\texttt{{-} Then a vertical line '│' (not a cell)} \\
\texttt{{-} The next cell is (4,4)? Actually, the coordinates are numbered 1 to 7, so after the line, the first cell is column 4? But the line is between columns.} \\
\\ 
\texttt{To avoid confusion, I should refer to the goal positions as originally agreed: O at (4,4). But (4,4) is occupied by a line. This means the proposed goal is invalid for column 4 because it has a vertical line.} \\
\\ 
\texttt{I need to communicate this to the other player. The goal configuration must avoid lines and borders.} \\
\\ 
\texttt{So I should send a message to adjust the goal.} \\
\\ 
\texttt{SAY: Wait, (4,4) has a vertical line '│' on my board, so it's not empty. We need to avoid lines. Let's shift all objects to row 5 instead? For example: A(1,5), D(2,5), M(3,5), O(4,5), P(5,5), S(6,5), T(7,5). What do you think?} \\
            }
        }
    }
     \\ \\

    \theutterance \stepcounter{utterance}  
    & & & \multicolumn{2}{p{0.3\linewidth}}{
        \cellcolor[rgb]{0.9,0.9,0.9}{
            \makecell[{{p{\linewidth}}}]{
                \texttt{\tiny{[GM$|$GM]}}
                \texttt{I see the issue: (4,4) is a vertical line '│', not an empty cell. The grid uses box{-}drawing characters, so I must avoid those. Looking at row 4: `║ADM│◌◌◌║ 4`, the cells are:} \\
\\ 
\texttt{{-} (1,4): A} \\
\texttt{{-} (2,4): D} \\
\texttt{{-} (3,4): M} \\
\texttt{{-} Then a vertical line '│' (not a cell)} \\
\texttt{{-} The next cell is (4,4)? Actually, the coordinates are numbered 1 to 7, so after the line, the first cell is column 4? But the line is between columns.} \\
\\ 
\texttt{To avoid confusion, I should refer to the goal positions as originally agreed: O at (4,4). But (4,4) is occupied by a line. This means the proposed goal is invalid for column 4 because it has a vertical line.} \\
\\ 
\texttt{I need to communicate this to the other player. The goal configuration must avoid lines and borders.} \\
\\ 
\texttt{So I should send a message to adjust the goal.} \\
\\ 
\texttt{SAY: Wait, (4,4) has a vertical line '│' on my board, so it's not empty. We need to avoid lines. Let's shift all objects to row 5 instead? For example: A(1,5), D(2,5), M(3,5), O(4,5), P(5,5), S(6,5), T(7,5). What do you think?} \\
            }
        }
    }
    & & \\ \\

    \theutterance \stepcounter{utterance}  
    & & & \multicolumn{2}{p{0.3\linewidth}}{
        \cellcolor[rgb]{0.9,0.9,0.9}{
            \makecell[{{p{\linewidth}}}]{
                \texttt{\tiny{[GM$|$GM]}}
                \texttt{Invalid format: head is not empty: 'I see the issue: (4,4) is a vertical line '│', not an empty cell. The grid uses box{-}drawing characters, so I must avoid those. Looking at row 4: ║ADM│◌◌◌║ 4, the cells are:} \\
\\ 
\texttt{{-} (1,4): A} \\
\texttt{{-} (2,4): D} \\
\texttt{{-} (3,4): M} \\
\texttt{{-} Then a vertical line '│' (not a cell)} \\
\texttt{{-} The next cell is (4,4)? Actually, the coordinates are numbered 1 to 7, so after the line, the first cell is column 4? But the line is between columns.} \\
\\ 
\texttt{To avoid confusion, I should refer to the goal positions as originally agreed: O at (4,4). But (4,4) is occupied by a line. This means the proposed goal is invalid for column 4 because it has a vertical line.} \\
\\ 
\texttt{I need to communicate this to the other player. The goal configuration must avoid lines and borders.} \\
\\ 
\texttt{So I should send a message to adjust the goal.} \\
\\ 
\texttt{'} \\
            }
        }
    }
    & & \\ \\

    \theutterance \stepcounter{utterance}  
    & & \multicolumn{4}{p{0.6\linewidth}}{
        \cellcolor[rgb]{0.9,0.9,0.9}{
            \makecell[{{p{\linewidth}}}]{
                \texttt{\tiny{[GM$\rangle$P2]}}
                \texttt{Penalty: Your message must not contain anything before the command!} \\
\\ 
\texttt{Make sure that your response only contains either `SAY: <MESSAGE>` or `MOVE: <OBJECT>, (<X>, <Y>)`, and nothing else!} \\
\\ 
\texttt{You have collectively accumulated 7 of 16 penalties.} \\
\texttt{Please try again!} \\
            }
        }
    }
    & \\ \\

    \theutterance \stepcounter{utterance}  
    & & & \multicolumn{4}{p{0.6\linewidth}}{
        \cellcolor[rgb]{0.9,0.9,0.9}{
            \makecell[{{p{\linewidth}}}]{
                \texttt{\tiny{[GM$\langle$P2]}}
                \texttt{SAY: Wait, (4,4) has a vertical line on my board. Let's shift to row 5: A(1,5), D(2,5), M(3,5), O(4,5), P(5,5), S(6,5), T(7,5). Agree?} \\
            }
        }
    }
     \\ \\

    \theutterance \stepcounter{utterance}  
    & & & \multicolumn{2}{p{0.3\linewidth}}{
        \cellcolor[rgb]{0.9,0.9,0.9}{
            \makecell[{{p{\linewidth}}}]{
                \texttt{\tiny{[GM$|$GM]}}
                \texttt{SAY: Wait, (4,4) has a vertical line on my board. Let's shift to row 5: A(1,5), D(2,5), M(3,5), O(4,5), P(5,5), S(6,5), T(7,5). Agree?} \\
            }
        }
    }
    & & \\ \\

    \theutterance \stepcounter{utterance}  
    & & \multicolumn{4}{p{0.6\linewidth}}{
        \cellcolor[rgb]{0.9,0.9,0.9}{
            \makecell[{{p{\linewidth}}}]{
                \texttt{\tiny{[P1$\langle$GM]}}
                \texttt{Your message has been relayed to the other player.} \\
\\ 
\\ 
\texttt{You are currently playing round 5 of maximum 28.} \\
\texttt{You have collectively accumulated 7 of 16 penalties.} \\
\texttt{The other player sent this message:} \\
\\ 
\\ 
\texttt{"Wait, (4,4) has a vertical line on my board. Let's shift to row 5: A(1,5), D(2,5), M(3,5), O(4,5), P(5,5), S(6,5), T(7,5). Agree?"} \\
\\ 
\\ 
\texttt{What is your next command?} \\
            }
        }
    }
    & \\ \\

    \theutterance \stepcounter{utterance}  
    & \multicolumn{4}{p{0.6\linewidth}}{
        \cellcolor[rgb]{0.9,0.9,0.9}{
            \makecell[{{p{\linewidth}}}]{
                \texttt{\tiny{[P1$\rangle$GM]}}
                \texttt{SAY: Yes, row 5 works better! I agree to new goal: A(1,5), D(2,5), M(3,5), O(4,5), P(5,5), S(6,5), T(7,5). I'll adjust my moves accordingly.} \\
            }
        }
    }
    & & \\ \\

    \theutterance \stepcounter{utterance}  
    & & & \multicolumn{2}{p{0.3\linewidth}}{
        \cellcolor[rgb]{0.9,0.9,0.9}{
            \makecell[{{p{\linewidth}}}]{
                \texttt{\tiny{[GM$|$GM]}}
                \texttt{SAY: Yes, row 5 works better! I agree to new goal: A(1,5), D(2,5), M(3,5), O(4,5), P(5,5), S(6,5), T(7,5). I'll adjust my moves accordingly.} \\
            }
        }
    }
    & & \\ \\

    \theutterance \stepcounter{utterance}  
    & & \multicolumn{4}{p{0.6\linewidth}}{
        \cellcolor[rgb]{0.9,0.9,0.9}{
            \makecell[{{p{\linewidth}}}]{
                \texttt{\tiny{[GM$\rangle$P2]}}
                \texttt{Your message has been relayed to the other player.} \\
\\ 
\\ 
\texttt{You are currently playing round 6 of maximum 28.} \\
\texttt{You have collectively accumulated 7 of 16 penalties.} \\
\texttt{The other player sent this message:} \\
\\ 
\\ 
\texttt{"Yes, row 5 works better! I agree to new goal: A(1,5), D(2,5), M(3,5), O(4,5), P(5,5), S(6,5), T(7,5). I'll adjust my moves accordingly."} \\
\\ 
\\ 
\texttt{What is your next command?} \\
            }
        }
    }
    & \\ \\

    \theutterance \stepcounter{utterance}  
    & & & \multicolumn{4}{p{0.6\linewidth}}{
        \cellcolor[rgb]{0.9,0.9,0.9}{
            \makecell[{{p{\linewidth}}}]{
                \texttt{\tiny{[GM$\langle$P2]}}
                \texttt{MOVE: A, (1,5)} \\
            }
        }
    }
     \\ \\

    \theutterance \stepcounter{utterance}  
    & & & \multicolumn{2}{p{0.3\linewidth}}{
        \cellcolor[rgb]{0.9,0.9,0.9}{
            \makecell[{{p{\linewidth}}}]{
                \texttt{\tiny{[GM$|$GM]}}
                \texttt{MOVE: A, (1,5)} \\
            }
        }
    }
    & & \\ \\

    \theutterance \stepcounter{utterance}  
    & & & \multicolumn{2}{p{0.3\linewidth}}{
        \cellcolor[rgb]{0.9,0.9,0.9}{
            \makecell[{{p{\linewidth}}}]{
                \texttt{\tiny{[GM$|$GM]}}
                \texttt{Penalty: (1,5) is not empty, but contains '─'.} \\
\\ 
\texttt{You have collectively accumulated 8 of 16 penalties.} \\
\texttt{Please try again!} \\
            }
        }
    }
    & & \\ \\

    \theutterance \stepcounter{utterance}  
    & & \multicolumn{4}{p{0.6\linewidth}}{
        \cellcolor[rgb]{0.9,0.9,0.9}{
            \makecell[{{p{\linewidth}}}]{
                \texttt{\tiny{[GM$\rangle$P2]}}
                \texttt{Penalty: (1,5) is not empty, but contains '─'.} \\
\\ 
\texttt{You have collectively accumulated 8 of 16 penalties.} \\
\texttt{Please try again!} \\
            }
        }
    }
    & \\ \\

    \theutterance \stepcounter{utterance}  
    & & & \multicolumn{4}{p{0.6\linewidth}}{
        \cellcolor[rgb]{0.9,0.9,0.9}{
            \makecell[{{p{\linewidth}}}]{
                \texttt{\tiny{[GM$\langle$P2]}}
                \texttt{Looking at row 5: `╟─┬─┼───╢ 5`, which is mostly lines and borders. So row 5 is not suitable for placing objects either. I need to choose a row that is entirely composed of empty cells. From the grid, row 4 is mostly empty except for the vertical line, but row 6 and row 7 have some objects and lines.} \\
\\ 
\texttt{Let's check row 3: `║◌◌◌├─┘◌║ 3` {-} this has lines and borders.} \\
\texttt{Row 2: `║◌O◌│◌├─╢ 2` {-} has lines.} \\
\texttt{Row 1: `║◌MS│◌│◌║ 1` {-} has lines.} \\
\\ 
\texttt{The only row that seems to have no lines is row 4, but it has a vertical line at column 4. However, the cells are at columns 1,2,3 and then 4,5,6,7? But the vertical line is at the boundary.} \\
\\ 
\texttt{Actually, the grid has 7 columns. The vertical lines are between columns 3 and 4, and between columns 5 and 6? This is messy.} \\
\\ 
\texttt{To avoid penalties, I should use a row that is shown to have only '◌' and no lines. Row 4 has `║◌◌◌│◌◌◌║`, which means columns 1,2,3 are empty and columns 4,5,6,7 are empty? But there is a line at column 4 boundary.} \\
\\ 
\texttt{Perhaps the cells are:} \\
\texttt{{-} Columns 1,2,3: empty} \\
\texttt{{-} Column 4: the line is not a cell, so column 4 is the first after the line? But the coordinate (4,4) would be that cell.} \\
\\ 
\texttt{In row 4: "║◌◌◌│◌◌◌║", the first three ◌ are columns 1,2,3. Then "│" is the separator between column 3 and 4. Then the next ◌ is column 4, then column 5 is ◌, column 6 is ◌, and column 7 is ◌? But there are only three ◌ after the line, so it might be columns 4,5,6, and column 7 is not shown? No, the right border is at the end.} \\
\\ 
\texttt{This is not working. I should use a different row. Row 6: `║◌│◌│◌◌D║ 6` {-} has lines.} \\
\texttt{Row 7: `║T│◌│◌PA║ 7` {-} has lines.} \\
\\ 
\texttt{The only row that might be fully available is row 4, but we must avoid the line cells. Since the line is at column 4, we should not use (4,4). So let's use only the cells that are ◌ in row 4: columns 1,2,3,5,6,7. But we have 7 objects, so we need 7 cells.} \\
\\ 
\texttt{This is not feasible. I should suggest using row 6 or row 7 after moving objects away.} \\
\\ 
\texttt{Given the penalties, I should communicate to choose a row that is clear.} \\
\\ 
\texttt{SAY: Row 5 has lines too. Let's use row 6: after moving current objects, we can use empty cells. Propose A(1,6), D(2,6), M(3,6), O(4,6), P(5,6), S(6,6), T(7,6). Agree?} \\
            }
        }
    }
     \\ \\

    \theutterance \stepcounter{utterance}  
    & & & \multicolumn{2}{p{0.3\linewidth}}{
        \cellcolor[rgb]{0.9,0.9,0.9}{
            \makecell[{{p{\linewidth}}}]{
                \texttt{\tiny{[GM$|$GM]}}
                \texttt{Looking at row 5: `╟─┬─┼───╢ 5`, which is mostly lines and borders. So row 5 is not suitable for placing objects either. I need to choose a row that is entirely composed of empty cells. From the grid, row 4 is mostly empty except for the vertical line, but row 6 and row 7 have some objects and lines.} \\
\\ 
\texttt{Let's check row 3: `║◌◌◌├─┘◌║ 3` {-} this has lines and borders.} \\
\texttt{Row 2: `║◌O◌│◌├─╢ 2` {-} has lines.} \\
\texttt{Row 1: `║◌MS│◌│◌║ 1` {-} has lines.} \\
\\ 
\texttt{The only row that seems to have no lines is row 4, but it has a vertical line at column 4. However, the cells are at columns 1,2,3 and then 4,5,6,7? But the vertical line is at the boundary.} \\
\\ 
\texttt{Actually, the grid has 7 columns. The vertical lines are between columns 3 and 4, and between columns 5 and 6? This is messy.} \\
\\ 
\texttt{To avoid penalties, I should use a row that is shown to have only '◌' and no lines. Row 4 has `║◌◌◌│◌◌◌║`, which means columns 1,2,3 are empty and columns 4,5,6,7 are empty? But there is a line at column 4 boundary.} \\
\\ 
\texttt{Perhaps the cells are:} \\
\texttt{{-} Columns 1,2,3: empty} \\
\texttt{{-} Column 4: the line is not a cell, so column 4 is the first after the line? But the coordinate (4,4) would be that cell.} \\
\\ 
\texttt{In row 4: "║◌◌◌│◌◌◌║", the first three ◌ are columns 1,2,3. Then "│" is the separator between column 3 and 4. Then the next ◌ is column 4, then column 5 is ◌, column 6 is ◌, and column 7 is ◌? But there are only three ◌ after the line, so it might be columns 4,5,6, and column 7 is not shown? No, the right border is at the end.} \\
\\ 
\texttt{This is not working. I should use a different row. Row 6: `║◌│◌│◌◌D║ 6` {-} has lines.} \\
\texttt{Row 7: `║T│◌│◌PA║ 7` {-} has lines.} \\
\\ 
\texttt{The only row that might be fully available is row 4, but we must avoid the line cells. Since the line is at column 4, we should not use (4,4). So let's use only the cells that are ◌ in row 4: columns 1,2,3,5,6,7. But we have 7 objects, so we need 7 cells.} \\
\\ 
\texttt{This is not feasible. I should suggest using row 6 or row 7 after moving objects away.} \\
\\ 
\texttt{Given the penalties, I should communicate to choose a row that is clear.} \\
\\ 
\texttt{SAY: Row 5 has lines too. Let's use row 6: after moving current objects, we can use empty cells. Propose A(1,6), D(2,6), M(3,6), O(4,6), P(5,6), S(6,6), T(7,6). Agree?} \\
            }
        }
    }
    & & \\ \\

    \theutterance \stepcounter{utterance}  
    & & & \multicolumn{2}{p{0.3\linewidth}}{
        \cellcolor[rgb]{0.9,0.9,0.9}{
            \makecell[{{p{\linewidth}}}]{
                \texttt{\tiny{[GM$|$GM]}}
                \texttt{Invalid format: head is not empty: 'Looking at row 5: ╟─┬─┼───╢ 5, which is mostly lines and borders. So row 5 is not suitable for placing objects either. I need to choose a row that is entirely composed of empty cells. From the grid, row 4 is mostly empty except for the vertical line, but row 6 and row 7 have some objects and lines.} \\
\\ 
\texttt{Let's check row 3: ║◌◌◌├─┘◌║ 3 {-} this has lines and borders.} \\
\texttt{Row 2: ║◌O◌│◌├─╢ 2 {-} has lines.} \\
\texttt{Row 1: ║◌MS│◌│◌║ 1 {-} has lines.} \\
\\ 
\texttt{The only row that seems to have no lines is row 4, but it has a vertical line at column 4. However, the cells are at columns 1,2,3 and then 4,5,6,7? But the vertical line is at the boundary.} \\
\\ 
\texttt{Actually, the grid has 7 columns. The vertical lines are between columns 3 and 4, and between columns 5 and 6? This is messy.} \\
\\ 
\texttt{To avoid penalties, I should use a row that is shown to have only '◌' and no lines. Row 4 has ║◌◌◌│◌◌◌║, which means columns 1,2,3 are empty and columns 4,5,6,7 are empty? But there is a line at column 4 boundary.} \\
\\ 
\texttt{Perhaps the cells are:} \\
\texttt{{-} Columns 1,2,3: empty} \\
\texttt{{-} Column 4: the line is not a cell, so column 4 is the first after the line? But the coordinate (4,4) would be that cell.} \\
\\ 
\texttt{In row 4: "║◌◌◌│◌◌◌║", the first three ◌ are columns 1,2,3. Then "│" is the separator between column 3 and 4. Then the next ◌ is column 4, then column 5 is ◌, column 6 is ◌, and column 7 is ◌? But there are only three ◌ after the line, so it might be columns 4,5,6, and column 7 is not shown? No, the right border is at the end.} \\
\\ 
\texttt{This is not working. I should use a different row. Row 6: ║◌│◌│◌◌D║ 6 {-} has lines.} \\
\texttt{Row 7: ║T│◌│◌PA║ 7 {-} has lines.} \\
\\ 
\texttt{The only row that might be fully available is row 4, but we must avoid the line cells. Since the line is at column 4, we should not use (4,4). So let's use only the cells that are ◌ in row 4: columns 1,2,3,5,6,7. But we have 7 objects, so we need 7 cells.} \\
\\ 
\texttt{This is not feasible. I should suggest using row 6 or row 7 after moving objects away.} \\
\\ 
\texttt{Given the penalties, I should communicate to choose a row that is clear.} \\
\\ 
\texttt{'} \\
            }
        }
    }
    & & \\ \\

    \theutterance \stepcounter{utterance}  
    & & \multicolumn{4}{p{0.6\linewidth}}{
        \cellcolor[rgb]{0.9,0.9,0.9}{
            \makecell[{{p{\linewidth}}}]{
                \texttt{\tiny{[GM$\rangle$P2]}}
                \texttt{Penalty: Your message must not contain anything before the command!} \\
\\ 
\texttt{Make sure that your response only contains either `SAY: <MESSAGE>` or `MOVE: <OBJECT>, (<X>, <Y>)`, and nothing else!} \\
\\ 
\texttt{You have collectively accumulated 9 of 16 penalties.} \\
\texttt{Please try again!} \\
            }
        }
    }
    & \\ \\

    \theutterance \stepcounter{utterance}  
    & & & \multicolumn{4}{p{0.6\linewidth}}{
        \cellcolor[rgb]{0.9,0.9,0.9}{
            \makecell[{{p{\linewidth}}}]{
                \texttt{\tiny{[GM$\langle$P2]}}
                \texttt{SAY: Row 5 has lines. Let's use row 6: A(1,6), D(2,6), M(3,6), O(4,6), P(5,6), S(6,6), T(7,6). Agree?} \\
            }
        }
    }
     \\ \\

    \theutterance \stepcounter{utterance}  
    & & & \multicolumn{2}{p{0.3\linewidth}}{
        \cellcolor[rgb]{0.9,0.9,0.9}{
            \makecell[{{p{\linewidth}}}]{
                \texttt{\tiny{[GM$|$GM]}}
                \texttt{SAY: Row 5 has lines. Let's use row 6: A(1,6), D(2,6), M(3,6), O(4,6), P(5,6), S(6,6), T(7,6). Agree?} \\
            }
        }
    }
    & & \\ \\

    \theutterance \stepcounter{utterance}  
    & & \multicolumn{4}{p{0.6\linewidth}}{
        \cellcolor[rgb]{0.9,0.9,0.9}{
            \makecell[{{p{\linewidth}}}]{
                \texttt{\tiny{[P1$\langle$GM]}}
                \texttt{Your message has been relayed to the other player.} \\
\\ 
\\ 
\texttt{You are currently playing round 6 of maximum 28.} \\
\texttt{You have collectively accumulated 9 of 16 penalties.} \\
\texttt{The other player sent this message:} \\
\\ 
\\ 
\texttt{"Row 5 has lines. Let's use row 6: A(1,6), D(2,6), M(3,6), O(4,6), P(5,6), S(6,6), T(7,6). Agree?"} \\
\\ 
\\ 
\texttt{What is your next command?} \\
            }
        }
    }
    & \\ \\

    \theutterance \stepcounter{utterance}  
    & \multicolumn{4}{p{0.6\linewidth}}{
        \cellcolor[rgb]{0.9,0.9,0.9}{
            \makecell[{{p{\linewidth}}}]{
                \texttt{\tiny{[P1$\rangle$GM]}}
                \texttt{SAY: Row 6 has objects too. Let's use row 3: A(1,3), D(2,3), M(3,3), O(4,3), P(5,3), S(6,3), T(7,3). That row is mostly empty. Agree?} \\
            }
        }
    }
    & & \\ \\

    \theutterance \stepcounter{utterance}  
    & & & \multicolumn{2}{p{0.3\linewidth}}{
        \cellcolor[rgb]{0.9,0.9,0.9}{
            \makecell[{{p{\linewidth}}}]{
                \texttt{\tiny{[GM$|$GM]}}
                \texttt{SAY: Row 6 has objects too. Let's use row 3: A(1,3), D(2,3), M(3,3), O(4,3), P(5,3), S(6,3), T(7,3). That row is mostly empty. Agree?} \\
            }
        }
    }
    & & \\ \\

    \theutterance \stepcounter{utterance}  
    & & \multicolumn{4}{p{0.6\linewidth}}{
        \cellcolor[rgb]{0.9,0.9,0.9}{
            \makecell[{{p{\linewidth}}}]{
                \texttt{\tiny{[GM$\rangle$P2]}}
                \texttt{Your message has been relayed to the other player.} \\
\\ 
\\ 
\texttt{You are currently playing round 7 of maximum 28.} \\
\texttt{You have collectively accumulated 9 of 16 penalties.} \\
\texttt{The other player sent this message:} \\
\\ 
\\ 
\texttt{"Row 6 has objects too. Let's use row 3: A(1,3), D(2,3), M(3,3), O(4,3), P(5,3), S(6,3), T(7,3). That row is mostly empty. Agree?"} \\
\\ 
\\ 
\texttt{What is your next command?} \\
            }
        }
    }
    & \\ \\

    \theutterance \stepcounter{utterance}  
    & & & \multicolumn{4}{p{0.6\linewidth}}{
        \cellcolor[rgb]{0.9,0.9,0.9}{
            \makecell[{{p{\linewidth}}}]{
                \texttt{\tiny{[GM$\langle$P2]}}
                \texttt{Looking at row 3: `║◌◌◌├─┘◌║ 3` {-} this row has lines and borders, so it's not entirely empty. However, the other player suggests using it, so perhaps on their board it is clear. To avoid further penalties, I should agree and proceed, but first, I need to check if the target cells are empty.} \\
\\ 
\texttt{In row 3:} \\
\texttt{{-} (1,3): ◌ (empty)} \\
\texttt{{-} (2,3): ◌ (empty)} \\
\texttt{{-} (3,3): ◌ (empty)} \\
\texttt{{-} (4,3): after "├", which is a line, not empty} \\
\texttt{{-} (5,3): "─", not empty} \\
\texttt{{-} (6,3): "┘", not empty} \\
\texttt{{-} (7,3): ◌ (empty)} \\
\\ 
\texttt{So only (1,3), (2,3), (3,3), and (7,3) are empty. We cannot use (4,3), (5,3), (6,3) because they contain lines.} \\
\\ 
\texttt{This won't work for 7 objects. I must communicate this.} \\
\\ 
\texttt{SAY: Row 3 has lines at (4,3), (5,3), (6,3). Only (1,3),(2,3),(3,3),(7,3) are empty. We need a row with 7 empty cells. Let's use row 4: we can use (1,4),(2,4),(3,4),(5,4),(6,4),(7,4) and avoid (4,4). But we have 7 objects, so we need one more. Perhaps row 5 has some empty cells?} \\
            }
        }
    }
     \\ \\

    \theutterance \stepcounter{utterance}  
    & & & \multicolumn{2}{p{0.3\linewidth}}{
        \cellcolor[rgb]{0.9,0.9,0.9}{
            \makecell[{{p{\linewidth}}}]{
                \texttt{\tiny{[GM$|$GM]}}
                \texttt{Looking at row 3: `║◌◌◌├─┘◌║ 3` {-} this row has lines and borders, so it's not entirely empty. However, the other player suggests using it, so perhaps on their board it is clear. To avoid further penalties, I should agree and proceed, but first, I need to check if the target cells are empty.} \\
\\ 
\texttt{In row 3:} \\
\texttt{{-} (1,3): ◌ (empty)} \\
\texttt{{-} (2,3): ◌ (empty)} \\
\texttt{{-} (3,3): ◌ (empty)} \\
\texttt{{-} (4,3): after "├", which is a line, not empty} \\
\texttt{{-} (5,3): "─", not empty} \\
\texttt{{-} (6,3): "┘", not empty} \\
\texttt{{-} (7,3): ◌ (empty)} \\
\\ 
\texttt{So only (1,3), (2,3), (3,3), and (7,3) are empty. We cannot use (4,3), (5,3), (6,3) because they contain lines.} \\
\\ 
\texttt{This won't work for 7 objects. I must communicate this.} \\
\\ 
\texttt{SAY: Row 3 has lines at (4,3), (5,3), (6,3). Only (1,3),(2,3),(3,3),(7,3) are empty. We need a row with 7 empty cells. Let's use row 4: we can use (1,4),(2,4),(3,4),(5,4),(6,4),(7,4) and avoid (4,4). But we have 7 objects, so we need one more. Perhaps row 5 has some empty cells?} \\
            }
        }
    }
    & & \\ \\

    \theutterance \stepcounter{utterance}  
    & & & \multicolumn{2}{p{0.3\linewidth}}{
        \cellcolor[rgb]{0.9,0.9,0.9}{
            \makecell[{{p{\linewidth}}}]{
                \texttt{\tiny{[GM$|$GM]}}
                \texttt{Invalid format: head is not empty: 'Looking at row 3: ║◌◌◌├─┘◌║ 3 {-} this row has lines and borders, so it's not entirely empty. However, the other player suggests using it, so perhaps on their board it is clear. To avoid further penalties, I should agree and proceed, but first, I need to check if the target cells are empty.} \\
\\ 
\texttt{In row 3:} \\
\texttt{{-} (1,3): ◌ (empty)} \\
\texttt{{-} (2,3): ◌ (empty)} \\
\texttt{{-} (3,3): ◌ (empty)} \\
\texttt{{-} (4,3): after "├", which is a line, not empty} \\
\texttt{{-} (5,3): "─", not empty} \\
\texttt{{-} (6,3): "┘", not empty} \\
\texttt{{-} (7,3): ◌ (empty)} \\
\\ 
\texttt{So only (1,3), (2,3), (3,3), and (7,3) are empty. We cannot use (4,3), (5,3), (6,3) because they contain lines.} \\
\\ 
\texttt{This won't work for 7 objects. I must communicate this.} \\
\\ 
\texttt{'} \\
            }
        }
    }
    & & \\ \\

    \theutterance \stepcounter{utterance}  
    & & \multicolumn{4}{p{0.6\linewidth}}{
        \cellcolor[rgb]{0.9,0.9,0.9}{
            \makecell[{{p{\linewidth}}}]{
                \texttt{\tiny{[GM$\rangle$P2]}}
                \texttt{Penalty: Your message must not contain anything before the command!} \\
\\ 
\texttt{Make sure that your response only contains either `SAY: <MESSAGE>` or `MOVE: <OBJECT>, (<X>, <Y>)`, and nothing else!} \\
\\ 
\texttt{You have collectively accumulated 10 of 16 penalties.} \\
\texttt{Please try again!} \\
            }
        }
    }
    & \\ \\

    \theutterance \stepcounter{utterance}  
    & & & \multicolumn{4}{p{0.6\linewidth}}{
        \cellcolor[rgb]{0.9,0.9,0.9}{
            \makecell[{{p{\linewidth}}}]{
                \texttt{\tiny{[GM$\langle$P2]}}
                \texttt{SAY: Row 3 has lines at (4,3),(5,3),(6,3). Only 4 empty cells. Row 4 has (1,4),(2,4),(3,4),(5,4),(6,4),(7,4) empty. We need 7 cells. Suggest using row 4 and row 5 combined?} \\
            }
        }
    }
     \\ \\

    \theutterance \stepcounter{utterance}  
    & & & \multicolumn{2}{p{0.3\linewidth}}{
        \cellcolor[rgb]{0.9,0.9,0.9}{
            \makecell[{{p{\linewidth}}}]{
                \texttt{\tiny{[GM$|$GM]}}
                \texttt{SAY: Row 3 has lines at (4,3),(5,3),(6,3). Only 4 empty cells. Row 4 has (1,4),(2,4),(3,4),(5,4),(6,4),(7,4) empty. We need 7 cells. Suggest using row 4 and row 5 combined?} \\
            }
        }
    }
    & & \\ \\

    \theutterance \stepcounter{utterance}  
    & & \multicolumn{4}{p{0.6\linewidth}}{
        \cellcolor[rgb]{0.9,0.9,0.9}{
            \makecell[{{p{\linewidth}}}]{
                \texttt{\tiny{[P1$\langle$GM]}}
                \texttt{Your message has been relayed to the other player.} \\
\\ 
\\ 
\texttt{You are currently playing round 7 of maximum 28.} \\
\texttt{You have collectively accumulated 10 of 16 penalties.} \\
\texttt{The other player sent this message:} \\
\\ 
\\ 
\texttt{"Row 3 has lines at (4,3),(5,3),(6,3). Only 4 empty cells. Row 4 has (1,4),(2,4),(3,4),(5,4),(6,4),(7,4) empty. We need 7 cells. Suggest using row 4 and row 5 combined?"} \\
\\ 
\\ 
\texttt{What is your next command?} \\
            }
        }
    }
    & \\ \\

    \theutterance \stepcounter{utterance}  
    & \multicolumn{4}{p{0.6\linewidth}}{
        \cellcolor[rgb]{0.9,0.9,0.9}{
            \makecell[{{p{\linewidth}}}]{
                \texttt{\tiny{[P1$\rangle$GM]}}
                \texttt{SAY: Let's use two rows to avoid lines. How about: A(1,4), D(2,4), M(3,4), O(4,4), P(5,4), S(6,4), T(7,4) for row 4 and if not possible, use row 5 for some? But row 4 has 6 empty cells (1,4),(2,4),(3,4),(5,4),(6,4),(7,4) and row 5 has (1,5),(2,5),(3,5),(5,5),(6,5),(7,5). We can split: put A,D,M in row 4 and O,P,S,T in row 5. Agree?} \\
            }
        }
    }
    & & \\ \\

    \theutterance \stepcounter{utterance}  
    & & & \multicolumn{2}{p{0.3\linewidth}}{
        \cellcolor[rgb]{0.9,0.9,0.9}{
            \makecell[{{p{\linewidth}}}]{
                \texttt{\tiny{[GM$|$GM]}}
                \texttt{SAY: Let's use two rows to avoid lines. How about: A(1,4), D(2,4), M(3,4), O(4,4), P(5,4), S(6,4), T(7,4) for row 4 and if not possible, use row 5 for some? But row 4 has 6 empty cells (1,4),(2,4),(3,4),(5,4),(6,4),(7,4) and row 5 has (1,5),(2,5),(3,5),(5,5),(6,5),(7,5). We can split: put A,D,M in row 4 and O,P,S,T in row 5. Agree?} \\
            }
        }
    }
    & & \\ \\

    \theutterance \stepcounter{utterance}  
    & & \multicolumn{4}{p{0.6\linewidth}}{
        \cellcolor[rgb]{0.9,0.9,0.9}{
            \makecell[{{p{\linewidth}}}]{
                \texttt{\tiny{[GM$\rangle$P2]}}
                \texttt{Your message has been relayed to the other player.} \\
\\ 
\\ 
\texttt{You are currently playing round 8 of maximum 28.} \\
\texttt{You have collectively accumulated 10 of 16 penalties.} \\
\texttt{The other player sent this message:} \\
\\ 
\\ 
\texttt{"Let's use two rows to avoid lines. How about: A(1,4), D(2,4), M(3,4), O(4,4), P(5,4), S(6,4), T(7,4) for row 4 and if not possible, use row 5 for some? But row 4 has 6 empty cells (1,4),(2,4),(3,4),(5,4),(6,4),(7,4) and row 5 has (1,5),(2,5),(3,5),(5,5),(6,5),(7,5). We can split: put A,D,M in row 4 and O,P,S,T in row 5. Agree?"} \\
\\ 
\\ 
\texttt{What is your next command?} \\
            }
        }
    }
    & \\ \\

    \theutterance \stepcounter{utterance}  
    & & & \multicolumn{4}{p{0.6\linewidth}}{
        \cellcolor[rgb]{0.9,0.9,0.9}{
            \makecell[{{p{\linewidth}}}]{
                \texttt{\tiny{[GM$\langle$P2]}}
                \texttt{SAY: Agree. Let's put A(1,4), D(2,4), M(3,4) in row 4 and O(1,5), P(2,5), S(3,5), T(5,5) in row 5. But we have 7 objects: A,D,M,O,P,S,T. So O,P,S,T in row 5: O(1,5), P(2,5), S(3,5), T(5,5). This uses 4 cells in row 5. Confirm?} \\
            }
        }
    }
     \\ \\

    \theutterance \stepcounter{utterance}  
    & & & \multicolumn{2}{p{0.3\linewidth}}{
        \cellcolor[rgb]{0.9,0.9,0.9}{
            \makecell[{{p{\linewidth}}}]{
                \texttt{\tiny{[GM$|$GM]}}
                \texttt{SAY: Agree. Let's put A(1,4), D(2,4), M(3,4) in row 4 and O(1,5), P(2,5), S(3,5), T(5,5) in row 5. But we have 7 objects: A,D,M,O,P,S,T. So O,P,S,T in row 5: O(1,5), P(2,5), S(3,5), T(5,5). This uses 4 cells in row 5. Confirm?} \\
            }
        }
    }
    & & \\ \\

    \theutterance \stepcounter{utterance}  
    & & \multicolumn{4}{p{0.6\linewidth}}{
        \cellcolor[rgb]{0.9,0.9,0.9}{
            \makecell[{{p{\linewidth}}}]{
                \texttt{\tiny{[P1$\langle$GM]}}
                \texttt{Your message has been relayed to the other player.} \\
\\ 
\\ 
\texttt{You are currently playing round 8 of maximum 28.} \\
\texttt{You have collectively accumulated 10 of 16 penalties.} \\
\texttt{The other player sent this message:} \\
\\ 
\\ 
\texttt{"Agree. Let's put A(1,4), D(2,4), M(3,4) in row 4 and O(1,5), P(2,5), S(3,5), T(5,5) in row 5. But we have 7 objects: A,D,M,O,P,S,T. So O,P,S,T in row 5: O(1,5), P(2,5), S(3,5), T(5,5). This uses 4 cells in row 5. Confirm?"} \\
\\ 
\\ 
\texttt{What is your next command?} \\
            }
        }
    }
    & \\ \\

    \theutterance \stepcounter{utterance}  
    & \multicolumn{4}{p{0.6\linewidth}}{
        \cellcolor[rgb]{0.9,0.9,0.9}{
            \makecell[{{p{\linewidth}}}]{
                \texttt{\tiny{[P1$\rangle$GM]}}
                \texttt{SAY: Confirmed! Goal: A(1,4), D(2,4), M(3,4) in row 4 and O(1,5), P(2,5), S(3,5), T(5,5) in row 5. I'll start moving.} \\
            }
        }
    }
    & & \\ \\

    \theutterance \stepcounter{utterance}  
    & & & \multicolumn{2}{p{0.3\linewidth}}{
        \cellcolor[rgb]{0.9,0.9,0.9}{
            \makecell[{{p{\linewidth}}}]{
                \texttt{\tiny{[GM$|$GM]}}
                \texttt{SAY: Confirmed! Goal: A(1,4), D(2,4), M(3,4) in row 4 and O(1,5), P(2,5), S(3,5), T(5,5) in row 5. I'll start moving.} \\
            }
        }
    }
    & & \\ \\

    \theutterance \stepcounter{utterance}  
    & & \multicolumn{4}{p{0.6\linewidth}}{
        \cellcolor[rgb]{0.9,0.9,0.9}{
            \makecell[{{p{\linewidth}}}]{
                \texttt{\tiny{[GM$\rangle$P2]}}
                \texttt{Your message has been relayed to the other player.} \\
\\ 
\\ 
\texttt{You are currently playing round 9 of maximum 28.} \\
\texttt{You have collectively accumulated 10 of 16 penalties.} \\
\texttt{The other player sent this message:} \\
\\ 
\\ 
\texttt{"Confirmed! Goal: A(1,4), D(2,4), M(3,4) in row 4 and O(1,5), P(2,5), S(3,5), T(5,5) in row 5. I'll start moving."} \\
\\ 
\\ 
\texttt{What is your next command?} \\
            }
        }
    }
    & \\ \\

    \theutterance \stepcounter{utterance}  
    & & & \multicolumn{4}{p{0.6\linewidth}}{
        \cellcolor[rgb]{0.9,0.9,0.9}{
            \makecell[{{p{\linewidth}}}]{
                \texttt{\tiny{[GM$\langle$P2]}}
                \texttt{MOVE: D, (2,4)} \\
            }
        }
    }
     \\ \\

    \theutterance \stepcounter{utterance}  
    & & & \multicolumn{2}{p{0.3\linewidth}}{
        \cellcolor[rgb]{0.9,0.9,0.9}{
            \makecell[{{p{\linewidth}}}]{
                \texttt{\tiny{[GM$|$GM]}}
                \texttt{MOVE: D, (2,4)} \\
            }
        }
    }
    & & \\ \\

    \theutterance \stepcounter{utterance}  
    & & & \multicolumn{2}{p{0.3\linewidth}}{
        \cellcolor[rgb]{0.9,0.9,0.9}{
            \makecell[{{p{\linewidth}}}]{
                \texttt{\tiny{[GM$|$GM]}}
                \texttt{Penalty: (2,4) is not empty, but contains 'D'.} \\
\\ 
\texttt{You have collectively accumulated 11 of 16 penalties.} \\
\texttt{Please try again!} \\
            }
        }
    }
    & & \\ \\

    \theutterance \stepcounter{utterance}  
    & & \multicolumn{4}{p{0.6\linewidth}}{
        \cellcolor[rgb]{0.9,0.9,0.9}{
            \makecell[{{p{\linewidth}}}]{
                \texttt{\tiny{[GM$\rangle$P2]}}
                \texttt{Penalty: (2,4) is not empty, but contains 'D'.} \\
\\ 
\texttt{You have collectively accumulated 11 of 16 penalties.} \\
\texttt{Please try again!} \\
            }
        }
    }
    & \\ \\

    \theutterance \stepcounter{utterance}  
    & & & \multicolumn{4}{p{0.6\linewidth}}{
        \cellcolor[rgb]{0.9,0.9,0.9}{
            \makecell[{{p{\linewidth}}}]{
                \texttt{\tiny{[GM$\langle$P2]}}
                \texttt{MOVE: O, (1,5)} \\
            }
        }
    }
     \\ \\

    \theutterance \stepcounter{utterance}  
    & & & \multicolumn{2}{p{0.3\linewidth}}{
        \cellcolor[rgb]{0.9,0.9,0.9}{
            \makecell[{{p{\linewidth}}}]{
                \texttt{\tiny{[GM$|$GM]}}
                \texttt{MOVE: O, (1,5)} \\
            }
        }
    }
    & & \\ \\

    \theutterance \stepcounter{utterance}  
    & & & \multicolumn{2}{p{0.3\linewidth}}{
        \cellcolor[rgb]{0.9,0.9,0.9}{
            \makecell[{{p{\linewidth}}}]{
                \texttt{\tiny{[GM$|$GM]}}
                \texttt{Penalty: (1,5) is not empty, but contains '─'.} \\
\\ 
\texttt{You have collectively accumulated 12 of 16 penalties.} \\
\texttt{Please try again!} \\
            }
        }
    }
    & & \\ \\

    \theutterance \stepcounter{utterance}  
    & & \multicolumn{4}{p{0.6\linewidth}}{
        \cellcolor[rgb]{0.9,0.9,0.9}{
            \makecell[{{p{\linewidth}}}]{
                \texttt{\tiny{[GM$\rangle$P2]}}
                \texttt{Penalty: (1,5) is not empty, but contains '─'.} \\
\\ 
\texttt{You have collectively accumulated 12 of 16 penalties.} \\
\texttt{Please try again!} \\
            }
        }
    }
    & \\ \\

    \theutterance \stepcounter{utterance}  
    & & & \multicolumn{4}{p{0.6\linewidth}}{
        \cellcolor[rgb]{0.9,0.9,0.9}{
            \makecell[{{p{\linewidth}}}]{
                \texttt{\tiny{[GM$\langle$P2]}}
                \texttt{Looking at my grid, I need to move objects to the agreed positions. However, (1,5) is not empty; it contains '─'. So I should avoid row 5 cells that have lines.} \\
\\ 
\texttt{The agreed goal for row 5 is O(1,5), P(2,5), S(3,5), T(5,5). But (1,5), (2,5), (3,5), and (5,5) might not be empty due to lines.} \\
\\ 
\texttt{I should check row 5 in my grid: `╟─┬─┼───╢ 5`} \\
\texttt{{-} (1,5): '─' (line)} \\
\texttt{{-} (2,5): '┬' (line)} \\
\texttt{{-} (3,5): '─' (line)} \\
\texttt{{-} (4,5): '┼' (line)} \\
\texttt{{-} (5,5): '─' (line)} \\
\texttt{{-} (6,5): '─' (line)} \\
\texttt{{-} (7,5): '─' (line)} \\
\\ 
\texttt{All cells in row 5 are lines! So I cannot place any object in row 5.} \\
\\ 
\texttt{I must communicate this to the other player.} \\
\\ 
\texttt{SAY: Row 5 is all lines: (1,5) to (7,5) are lines. Cannot use row 5. Let's use row 6 instead: O(1,6), P(2,6), S(3,6), T(5,6). But check if empty.} \\
            }
        }
    }
     \\ \\

    \theutterance \stepcounter{utterance}  
    & & & \multicolumn{2}{p{0.3\linewidth}}{
        \cellcolor[rgb]{0.9,0.9,0.9}{
            \makecell[{{p{\linewidth}}}]{
                \texttt{\tiny{[GM$|$GM]}}
                \texttt{Looking at my grid, I need to move objects to the agreed positions. However, (1,5) is not empty; it contains '─'. So I should avoid row 5 cells that have lines.} \\
\\ 
\texttt{The agreed goal for row 5 is O(1,5), P(2,5), S(3,5), T(5,5). But (1,5), (2,5), (3,5), and (5,5) might not be empty due to lines.} \\
\\ 
\texttt{I should check row 5 in my grid: `╟─┬─┼───╢ 5`} \\
\texttt{{-} (1,5): '─' (line)} \\
\texttt{{-} (2,5): '┬' (line)} \\
\texttt{{-} (3,5): '─' (line)} \\
\texttt{{-} (4,5): '┼' (line)} \\
\texttt{{-} (5,5): '─' (line)} \\
\texttt{{-} (6,5): '─' (line)} \\
\texttt{{-} (7,5): '─' (line)} \\
\\ 
\texttt{All cells in row 5 are lines! So I cannot place any object in row 5.} \\
\\ 
\texttt{I must communicate this to the other player.} \\
\\ 
\texttt{SAY: Row 5 is all lines: (1,5) to (7,5) are lines. Cannot use row 5. Let's use row 6 instead: O(1,6), P(2,6), S(3,6), T(5,6). But check if empty.} \\
            }
        }
    }
    & & \\ \\

    \theutterance \stepcounter{utterance}  
    & & & \multicolumn{2}{p{0.3\linewidth}}{
        \cellcolor[rgb]{0.9,0.9,0.9}{
            \makecell[{{p{\linewidth}}}]{
                \texttt{\tiny{[GM$|$GM]}}
                \texttt{Invalid format: head is not empty: 'Looking at my grid, I need to move objects to the agreed positions. However, (1,5) is not empty; it contains '─'. So I should avoid row 5 cells that have lines.} \\
\\ 
\texttt{The agreed goal for row 5 is O(1,5), P(2,5), S(3,5), T(5,5). But (1,5), (2,5), (3,5), and (5,5) might not be empty due to lines.} \\
\\ 
\texttt{I should check row 5 in my grid: ╟─┬─┼───╢ 5} \\
\texttt{{-} (1,5): '─' (line)} \\
\texttt{{-} (2,5): '┬' (line)} \\
\texttt{{-} (3,5): '─' (line)} \\
\texttt{{-} (4,5): '┼' (line)} \\
\texttt{{-} (5,5): '─' (line)} \\
\texttt{{-} (6,5): '─' (line)} \\
\texttt{{-} (7,5): '─' (line)} \\
\\ 
\texttt{All cells in row 5 are lines! So I cannot place any object in row 5.} \\
\\ 
\texttt{I must communicate this to the other player.} \\
\\ 
\texttt{'} \\
            }
        }
    }
    & & \\ \\

    \theutterance \stepcounter{utterance}  
    & & \multicolumn{4}{p{0.6\linewidth}}{
        \cellcolor[rgb]{0.9,0.9,0.9}{
            \makecell[{{p{\linewidth}}}]{
                \texttt{\tiny{[GM$\rangle$P2]}}
                \texttt{Penalty: Your message must not contain anything before the command!} \\
\\ 
\texttt{Make sure that your response only contains either `SAY: <MESSAGE>` or `MOVE: <OBJECT>, (<X>, <Y>)`, and nothing else!} \\
\\ 
\texttt{You have collectively accumulated 13 of 16 penalties.} \\
\texttt{Please try again!} \\
            }
        }
    }
    & \\ \\

    \theutterance \stepcounter{utterance}  
    & & & \multicolumn{4}{p{0.6\linewidth}}{
        \cellcolor[rgb]{0.9,0.9,0.9}{
            \makecell[{{p{\linewidth}}}]{
                \texttt{\tiny{[GM$\langle$P2]}}
                \texttt{SAY: Row 5 is all lines. Cannot use. Let's use row 6: O(1,6), P(2,6), S(3,6), T(5,6). Check if empty.} \\
            }
        }
    }
     \\ \\

    \theutterance \stepcounter{utterance}  
    & & & \multicolumn{2}{p{0.3\linewidth}}{
        \cellcolor[rgb]{0.9,0.9,0.9}{
            \makecell[{{p{\linewidth}}}]{
                \texttt{\tiny{[GM$|$GM]}}
                \texttt{SAY: Row 5 is all lines. Cannot use. Let's use row 6: O(1,6), P(2,6), S(3,6), T(5,6). Check if empty.} \\
            }
        }
    }
    & & \\ \\

    \theutterance \stepcounter{utterance}  
    & & \multicolumn{4}{p{0.6\linewidth}}{
        \cellcolor[rgb]{0.9,0.9,0.9}{
            \makecell[{{p{\linewidth}}}]{
                \texttt{\tiny{[P1$\langle$GM]}}
                \texttt{Your message has been relayed to the other player.} \\
\\ 
\\ 
\texttt{You are currently playing round 9 of maximum 28.} \\
\texttt{You have collectively accumulated 13 of 16 penalties.} \\
\texttt{The other player sent this message:} \\
\\ 
\\ 
\texttt{"Row 5 is all lines. Cannot use. Let's use row 6: O(1,6), P(2,6), S(3,6), T(5,6). Check if empty."} \\
\\ 
\\ 
\texttt{What is your next command?} \\
            }
        }
    }
    & \\ \\

    \theutterance \stepcounter{utterance}  
    & \multicolumn{4}{p{0.6\linewidth}}{
        \cellcolor[rgb]{0.9,0.9,0.9}{
            \makecell[{{p{\linewidth}}}]{
                \texttt{\tiny{[P1$\rangle$GM]}}
                \texttt{SAY: On my board, row 6 has S at (6,6) and O at (7,6) currently. So (1,6) is empty, (2,6) is empty, (3,6) is empty, (5,6) is empty. That works for O,P,S,T. So row 4: A(1,4), D(2,4), M(3,4) and row 6: O(1,6), P(2,6), S(3,6), T(5,6). Agree? I'll proceed.} \\
            }
        }
    }
    & & \\ \\

    \theutterance \stepcounter{utterance}  
    & & & \multicolumn{2}{p{0.3\linewidth}}{
        \cellcolor[rgb]{0.9,0.9,0.9}{
            \makecell[{{p{\linewidth}}}]{
                \texttt{\tiny{[GM$|$GM]}}
                \texttt{SAY: On my board, row 6 has S at (6,6) and O at (7,6) currently. So (1,6) is empty, (2,6) is empty, (3,6) is empty, (5,6) is empty. That works for O,P,S,T. So row 4: A(1,4), D(2,4), M(3,4) and row 6: O(1,6), P(2,6), S(3,6), T(5,6). Agree? I'll proceed.} \\
            }
        }
    }
    & & \\ \\

    \theutterance \stepcounter{utterance}  
    & & \multicolumn{4}{p{0.6\linewidth}}{
        \cellcolor[rgb]{0.9,0.9,0.9}{
            \makecell[{{p{\linewidth}}}]{
                \texttt{\tiny{[GM$\rangle$P2]}}
                \texttt{Your message has been relayed to the other player.} \\
\\ 
\\ 
\texttt{You are currently playing round 10 of maximum 28.} \\
\texttt{You have collectively accumulated 13 of 16 penalties.} \\
\texttt{The other player sent this message:} \\
\\ 
\\ 
\texttt{"On my board, row 6 has S at (6,6) and O at (7,6) currently. So (1,6) is empty, (2,6) is empty, (3,6) is empty, (5,6) is empty. That works for O,P,S,T. So row 4: A(1,4), D(2,4), M(3,4) and row 6: O(1,6), P(2,6), S(3,6), T(5,6). Agree? I'll proceed."} \\
\\ 
\\ 
\texttt{What is your next command?} \\
            }
        }
    }
    & \\ \\

    \theutterance \stepcounter{utterance}  
    & & & \multicolumn{4}{p{0.6\linewidth}}{
        \cellcolor[rgb]{0.9,0.9,0.9}{
            \makecell[{{p{\linewidth}}}]{
                \texttt{\tiny{[GM$\langle$P2]}}
                \texttt{MOVE: S, (3,6)} \\
            }
        }
    }
     \\ \\

    \theutterance \stepcounter{utterance}  
    & & & \multicolumn{2}{p{0.3\linewidth}}{
        \cellcolor[rgb]{0.9,0.9,0.9}{
            \makecell[{{p{\linewidth}}}]{
                \texttt{\tiny{[GM$|$GM]}}
                \texttt{MOVE: S, (3,6)} \\
            }
        }
    }
    & & \\ \\

    \theutterance \stepcounter{utterance}  
    & & & \multicolumn{2}{p{0.3\linewidth}}{
        \cellcolor[rgb]{0.9,0.9,0.9}{
            \makecell[{{p{\linewidth}}}]{
                \texttt{\tiny{[GM$|$GM]}}
                \texttt{Moved 'S' to (3,6) successfully. Your updated grid looks like this:} \\
\texttt{```} \\
\texttt{ 1234567} \\
\texttt{╔═══╤═╤═╗} \\
\texttt{║◌◌◌│◌│◌║ 1} \\
\texttt{║◌O◌│◌├─╢ 2} \\
\texttt{║◌◌◌├─┘◌║ 3} \\
\texttt{║ADM│◌◌◌║ 4} \\
\texttt{╟─┬─┼───╢ 5} \\
\texttt{║◌│S│◌◌◌║ 6} \\
\texttt{║T│◌│◌P◌║ 7} \\
\texttt{╚═╧═╧═══╝} \\
\texttt{```} \\
            }
        }
    }
    & & \\ \\

    \theutterance \stepcounter{utterance}  
    & & \multicolumn{4}{p{0.6\linewidth}}{
        \cellcolor[rgb]{0.9,0.9,0.9}{
            \makecell[{{p{\linewidth}}}]{
                \texttt{\tiny{[P1$\langle$GM]}}
                \texttt{Your message has been relayed to the other player.} \\
\\ 
\\ 
\texttt{You are currently playing round 10 of maximum 28.} \\
\texttt{You have collectively accumulated 13 of 16 penalties.} \\
\texttt{The other player moved an object on their grid.} \\
\\ 
\\ 
\texttt{What is your next command?} \\
            }
        }
    }
    & \\ \\

    \theutterance \stepcounter{utterance}  
    & \multicolumn{4}{p{0.6\linewidth}}{
        \cellcolor[rgb]{0.9,0.9,0.9}{
            \makecell[{{p{\linewidth}}}]{
                \texttt{\tiny{[P1$\rangle$GM]}}
                \texttt{MOVE: A, (1,4)} \\
            }
        }
    }
    & & \\ \\

    \theutterance \stepcounter{utterance}  
    & & & \multicolumn{2}{p{0.3\linewidth}}{
        \cellcolor[rgb]{0.9,0.9,0.9}{
            \makecell[{{p{\linewidth}}}]{
                \texttt{\tiny{[GM$|$GM]}}
                \texttt{MOVE: A, (1,4)} \\
            }
        }
    }
    & & \\ \\

    \theutterance \stepcounter{utterance}  
    & & & \multicolumn{2}{p{0.3\linewidth}}{
        \cellcolor[rgb]{0.9,0.9,0.9}{
            \makecell[{{p{\linewidth}}}]{
                \texttt{\tiny{[GM$|$GM]}}
                \texttt{Penalty: (1,4) is not empty, but contains 'A'.} \\
\\ 
\texttt{You have collectively accumulated 14 of 16 penalties.} \\
\texttt{Please try again!} \\
            }
        }
    }
    & & \\ \\

    \theutterance \stepcounter{utterance}  
    & & \multicolumn{4}{p{0.6\linewidth}}{
        \cellcolor[rgb]{0.9,0.9,0.9}{
            \makecell[{{p{\linewidth}}}]{
                \texttt{\tiny{[P1$\langle$GM]}}
                \texttt{Penalty: (1,4) is not empty, but contains 'A'.} \\
\\ 
\texttt{You have collectively accumulated 14 of 16 penalties.} \\
\texttt{Please try again!} \\
            }
        }
    }
    & \\ \\

    \theutterance \stepcounter{utterance}  
    & \multicolumn{4}{p{0.6\linewidth}}{
        \cellcolor[rgb]{0.9,0.9,0.9}{
            \makecell[{{p{\linewidth}}}]{
                \texttt{\tiny{[P1$\rangle$GM]}}
                \texttt{SAY: A is already at (1,4) on my board. I'll move D to (2,4).} \\
            }
        }
    }
    & & \\ \\

    \theutterance \stepcounter{utterance}  
    & & & \multicolumn{2}{p{0.3\linewidth}}{
        \cellcolor[rgb]{0.9,0.9,0.9}{
            \makecell[{{p{\linewidth}}}]{
                \texttt{\tiny{[GM$|$GM]}}
                \texttt{SAY: A is already at (1,4) on my board. I'll move D to (2,4).} \\
            }
        }
    }
    & & \\ \\

    \theutterance \stepcounter{utterance}  
    & & \multicolumn{4}{p{0.6\linewidth}}{
        \cellcolor[rgb]{0.9,0.9,0.9}{
            \makecell[{{p{\linewidth}}}]{
                \texttt{\tiny{[GM$\rangle$P2]}}
                \texttt{Moved 'S' to (3,6) successfully. Your updated grid looks like this:} \\
\texttt{```} \\
\texttt{ 1234567} \\
\texttt{╔═══╤═╤═╗} \\
\texttt{║◌◌◌│◌│◌║ 1} \\
\texttt{║◌O◌│◌├─╢ 2} \\
\texttt{║◌◌◌├─┘◌║ 3} \\
\texttt{║ADM│◌◌◌║ 4} \\
\texttt{╟─┬─┼───╢ 5} \\
\texttt{║◌│S│◌◌◌║ 6} \\
\texttt{║T│◌│◌P◌║ 7} \\
\texttt{╚═╧═╧═══╝} \\
\texttt{```} \\
\texttt{You are currently playing round 11 of maximum 28.} \\
\texttt{You have collectively accumulated 14 of 16 penalties.} \\
\texttt{The other player sent this message:} \\
\\ 
\\ 
\texttt{"A is already at (1,4) on my board. I'll move D to (2,4)."} \\
\\ 
\\ 
\texttt{What is your next command?} \\
            }
        }
    }
    & \\ \\

    \theutterance \stepcounter{utterance}  
    & & & \multicolumn{4}{p{0.6\linewidth}}{
        \cellcolor[rgb]{0.9,0.9,0.9}{
            \makecell[{{p{\linewidth}}}]{
                \texttt{\tiny{[GM$\langle$P2]}}
                \texttt{MOVE: P, (2,6)} \\
            }
        }
    }
     \\ \\

    \theutterance \stepcounter{utterance}  
    & & & \multicolumn{2}{p{0.3\linewidth}}{
        \cellcolor[rgb]{0.9,0.9,0.9}{
            \makecell[{{p{\linewidth}}}]{
                \texttt{\tiny{[GM$|$GM]}}
                \texttt{MOVE: P, (2,6)} \\
            }
        }
    }
    & & \\ \\

    \theutterance \stepcounter{utterance}  
    & & & \multicolumn{2}{p{0.3\linewidth}}{
        \cellcolor[rgb]{0.9,0.9,0.9}{
            \makecell[{{p{\linewidth}}}]{
                \texttt{\tiny{[GM$|$GM]}}
                \texttt{Penalty: (2,6) is not empty, but contains '│'.} \\
\\ 
\texttt{You have collectively accumulated 15 of 16 penalties.} \\
\texttt{Please try again!} \\
            }
        }
    }
    & & \\ \\

    \theutterance \stepcounter{utterance}  
    & & \multicolumn{4}{p{0.6\linewidth}}{
        \cellcolor[rgb]{0.9,0.9,0.9}{
            \makecell[{{p{\linewidth}}}]{
                \texttt{\tiny{[GM$\rangle$P2]}}
                \texttt{Penalty: (2,6) is not empty, but contains '│'.} \\
\\ 
\texttt{You have collectively accumulated 15 of 16 penalties.} \\
\texttt{Please try again!} \\
            }
        }
    }
    & \\ \\

    \theutterance \stepcounter{utterance}  
    & & & \multicolumn{4}{p{0.6\linewidth}}{
        \cellcolor[rgb]{0.9,0.9,0.9}{
            \makecell[{{p{\linewidth}}}]{
                \texttt{\tiny{[GM$\langle$P2]}}
                \texttt{Looking at row 6: `║◌│S│◌◌◌║ 6` {-} (2,6) is a vertical line '│', not empty. So I cannot move P to (2,6). I need to choose an empty cell in row 6.} \\
\\ 
\texttt{From the grid:} \\
\texttt{{-} (1,6): ◌ (empty)} \\
\texttt{{-} (2,6): │ (line)} \\
\texttt{{-} (3,6): S (now occupied)} \\
\texttt{{-} (4,6): │ (line)} \\
\texttt{{-} (5,6): ◌ (empty)} \\
\texttt{{-} (6,6): ◌ (empty)} \\
\texttt{{-} (7,6): ◌ (empty)} \\
\\ 
\texttt{So the empty cells in row 6 are (1,6), (5,6), (6,6), (7,6).} \\
\\ 
\texttt{The agreed goal for row 6 is O(1,6), P(2,6), S(3,6), T(5,6). But (2,6) is not available. I should adjust and use (1,6), (5,6), (6,6), (7,6) for O,P,S,T.} \\
\\ 
\texttt{I need to communicate this to the other player.} \\
\\ 
\texttt{SAY: (2,6) is a line on my board. Use (1,6) for O, (5,6) for T, (6,6) for P, (7,6) for S? But S is at (3,6). Let's do: O(1,6), P(6,6), S(7,6), T(5,6). Agree?} \\
            }
        }
    }
     \\ \\

    \theutterance \stepcounter{utterance}  
    & & & \multicolumn{2}{p{0.3\linewidth}}{
        \cellcolor[rgb]{0.9,0.9,0.9}{
            \makecell[{{p{\linewidth}}}]{
                \texttt{\tiny{[GM$|$GM]}}
                \texttt{Looking at row 6: `║◌│S│◌◌◌║ 6` {-} (2,6) is a vertical line '│', not empty. So I cannot move P to (2,6). I need to choose an empty cell in row 6.} \\
\\ 
\texttt{From the grid:} \\
\texttt{{-} (1,6): ◌ (empty)} \\
\texttt{{-} (2,6): │ (line)} \\
\texttt{{-} (3,6): S (now occupied)} \\
\texttt{{-} (4,6): │ (line)} \\
\texttt{{-} (5,6): ◌ (empty)} \\
\texttt{{-} (6,6): ◌ (empty)} \\
\texttt{{-} (7,6): ◌ (empty)} \\
\\ 
\texttt{So the empty cells in row 6 are (1,6), (5,6), (6,6), (7,6).} \\
\\ 
\texttt{The agreed goal for row 6 is O(1,6), P(2,6), S(3,6), T(5,6). But (2,6) is not available. I should adjust and use (1,6), (5,6), (6,6), (7,6) for O,P,S,T.} \\
\\ 
\texttt{I need to communicate this to the other player.} \\
\\ 
\texttt{SAY: (2,6) is a line on my board. Use (1,6) for O, (5,6) for T, (6,6) for P, (7,6) for S? But S is at (3,6). Let's do: O(1,6), P(6,6), S(7,6), T(5,6). Agree?} \\
            }
        }
    }
    & & \\ \\

    \theutterance \stepcounter{utterance}  
    & & & \multicolumn{2}{p{0.3\linewidth}}{
        \cellcolor[rgb]{0.9,0.9,0.9}{
            \makecell[{{p{\linewidth}}}]{
                \texttt{\tiny{[GM$|$GM]}}
                \texttt{Invalid format: head is not empty: 'Looking at row 6: ║◌│S│◌◌◌║ 6 {-} (2,6) is a vertical line '│', not empty. So I cannot move P to (2,6). I need to choose an empty cell in row 6.} \\
\\ 
\texttt{From the grid:} \\
\texttt{{-} (1,6): ◌ (empty)} \\
\texttt{{-} (2,6): │ (line)} \\
\texttt{{-} (3,6): S (now occupied)} \\
\texttt{{-} (4,6): │ (line)} \\
\texttt{{-} (5,6): ◌ (empty)} \\
\texttt{{-} (6,6): ◌ (empty)} \\
\texttt{{-} (7,6): ◌ (empty)} \\
\\ 
\texttt{So the empty cells in row 6 are (1,6), (5,6), (6,6), (7,6).} \\
\\ 
\texttt{The agreed goal for row 6 is O(1,6), P(2,6), S(3,6), T(5,6). But (2,6) is not available. I should adjust and use (1,6), (5,6), (6,6), (7,6) for O,P,S,T.} \\
\\ 
\texttt{I need to communicate this to the other player.} \\
\\ 
\texttt{'} \\
            }
        }
    }
    & & \\ \\

    \theutterance \stepcounter{utterance}  
    & & \multicolumn{4}{p{0.6\linewidth}}{
        \cellcolor[rgb]{0.9,0.9,0.9}{
            \makecell[{{p{\linewidth}}}]{
                \texttt{\tiny{[GM$\rangle$P2]}}
                \texttt{Penalty: Your message must not contain anything before the command!} \\
\\ 
\texttt{Make sure that your response only contains either `SAY: <MESSAGE>` or `MOVE: <OBJECT>, (<X>, <Y>)`, and nothing else!} \\
\\ 
\texttt{You have collectively accumulated 16 of 16 penalties.} \\
\texttt{Please try again!} \\
            }
        }
    }
    & \\ \\

    \theutterance \stepcounter{utterance}  
    & & & \multicolumn{4}{p{0.6\linewidth}}{
        \cellcolor[rgb]{0.9,0.9,0.9}{
            \makecell[{{p{\linewidth}}}]{
                \texttt{\tiny{[GM$\langle$P2]}}
                \texttt{SAY: (2,6) is line. Use O(1,6), P(6,6), S(7,6), T(5,6). Agree?} \\
            }
        }
    }
     \\ \\

    \theutterance \stepcounter{utterance}  
    & & & \multicolumn{2}{p{0.3\linewidth}}{
        \cellcolor[rgb]{0.9,0.9,0.9}{
            \makecell[{{p{\linewidth}}}]{
                \texttt{\tiny{[GM$|$GM]}}
                \texttt{SAY: (2,6) is line. Use O(1,6), P(6,6), S(7,6), T(5,6). Agree?} \\
            }
        }
    }
    & & \\ \\

    \theutterance \stepcounter{utterance}  
    & & \multicolumn{4}{p{0.6\linewidth}}{
        \cellcolor[rgb]{0.9,0.9,0.9}{
            \makecell[{{p{\linewidth}}}]{
                \texttt{\tiny{[P1$\langle$GM]}}
                \texttt{Your message has been relayed to the other player.} \\
\\ 
\\ 
\texttt{You are currently playing round 11 of maximum 28.} \\
\texttt{You have collectively accumulated 16 of 16 penalties.} \\
\texttt{The other player sent this message:} \\
\\ 
\\ 
\texttt{"(2,6) is line. Use O(1,6), P(6,6), S(7,6), T(5,6). Agree?"} \\
\\ 
\\ 
\texttt{What is your next command?} \\
            }
        }
    }
    & \\ \\

    \theutterance \stepcounter{utterance}  
    & \multicolumn{4}{p{0.6\linewidth}}{
        \cellcolor[rgb]{0.9,0.9,0.9}{
            \makecell[{{p{\linewidth}}}]{
                \texttt{\tiny{[P1$\rangle$GM]}}
                \texttt{SAY: Agreed! New row 6: O(1,6), P(6,6), S(7,6), T(5,6). And row 4: A(1,4), D(2,4), M(3,4). I'll move D to (2,4).} \\
            }
        }
    }
    & & \\ \\

    \theutterance \stepcounter{utterance}  
    & & & \multicolumn{2}{p{0.3\linewidth}}{
        \cellcolor[rgb]{0.9,0.9,0.9}{
            \makecell[{{p{\linewidth}}}]{
                \texttt{\tiny{[GM$|$GM]}}
                \texttt{SAY: Agreed! New row 6: O(1,6), P(6,6), S(7,6), T(5,6). And row 4: A(1,4), D(2,4), M(3,4). I'll move D to (2,4).} \\
            }
        }
    }
    & & \\ \\

    \theutterance \stepcounter{utterance}  
    & & \multicolumn{4}{p{0.6\linewidth}}{
        \cellcolor[rgb]{0.9,0.9,0.9}{
            \makecell[{{p{\linewidth}}}]{
                \texttt{\tiny{[GM$\rangle$P2]}}
                \texttt{Your message has been relayed to the other player.} \\
\\ 
\\ 
\texttt{You are currently playing round 12 of maximum 28.} \\
\texttt{You have collectively accumulated 16 of 16 penalties.} \\
\texttt{The other player sent this message:} \\
\\ 
\\ 
\texttt{"Agreed! New row 6: O(1,6), P(6,6), S(7,6), T(5,6). And row 4: A(1,4), D(2,4), M(3,4). I'll move D to (2,4)."} \\
\\ 
\\ 
\texttt{What is your next command?} \\
            }
        }
    }
    & \\ \\

    \theutterance \stepcounter{utterance}  
    & & & \multicolumn{4}{p{0.6\linewidth}}{
        \cellcolor[rgb]{0.9,0.9,0.9}{
            \makecell[{{p{\linewidth}}}]{
                \texttt{\tiny{[GM$\langle$P2]}}
                \texttt{MOVE: P, (6,6)} \\
            }
        }
    }
     \\ \\

    \theutterance \stepcounter{utterance}  
    & & & \multicolumn{2}{p{0.3\linewidth}}{
        \cellcolor[rgb]{0.9,0.9,0.9}{
            \makecell[{{p{\linewidth}}}]{
                \texttt{\tiny{[GM$|$GM]}}
                \texttt{MOVE: P, (6,6)} \\
            }
        }
    }
    & & \\ \\

    \theutterance \stepcounter{utterance}  
    & & & \multicolumn{2}{p{0.3\linewidth}}{
        \cellcolor[rgb]{0.9,0.9,0.9}{
            \makecell[{{p{\linewidth}}}]{
                \texttt{\tiny{[GM$|$GM]}}
                \texttt{Moved 'P' to (6,6) successfully. Your updated grid looks like this:} \\
\texttt{```} \\
\texttt{ 1234567} \\
\texttt{╔═══╤═╤═╗} \\
\texttt{║◌◌◌│◌│◌║ 1} \\
\texttt{║◌O◌│◌├─╢ 2} \\
\texttt{║◌◌◌├─┘◌║ 3} \\
\texttt{║ADM│◌◌◌║ 4} \\
\texttt{╟─┬─┼───╢ 5} \\
\texttt{║◌│S│◌P◌║ 6} \\
\texttt{║T│◌│◌◌◌║ 7} \\
\texttt{╚═╧═╧═══╝} \\
\texttt{```} \\
            }
        }
    }
    & & \\ \\

    \theutterance \stepcounter{utterance}  
    & & \multicolumn{4}{p{0.6\linewidth}}{
        \cellcolor[rgb]{0.9,0.9,0.9}{
            \makecell[{{p{\linewidth}}}]{
                \texttt{\tiny{[P1$\langle$GM]}}
                \texttt{Your message has been relayed to the other player.} \\
\\ 
\\ 
\texttt{You are currently playing round 12 of maximum 28.} \\
\texttt{You have collectively accumulated 16 of 16 penalties.} \\
\texttt{The other player moved an object on their grid.} \\
\\ 
\\ 
\texttt{What is your next command?} \\
            }
        }
    }
    & \\ \\

    \theutterance \stepcounter{utterance}  
    & \multicolumn{4}{p{0.6\linewidth}}{
        \cellcolor[rgb]{0.9,0.9,0.9}{
            \makecell[{{p{\linewidth}}}]{
                \texttt{\tiny{[P1$\rangle$GM]}}
                \texttt{MOVE: D, (2,4)} \\
            }
        }
    }
    & & \\ \\

    \theutterance \stepcounter{utterance}  
    & & & \multicolumn{2}{p{0.3\linewidth}}{
        \cellcolor[rgb]{0.9,0.9,0.9}{
            \makecell[{{p{\linewidth}}}]{
                \texttt{\tiny{[GM$|$GM]}}
                \texttt{MOVE: D, (2,4)} \\
            }
        }
    }
    & & \\ \\

    \theutterance \stepcounter{utterance}  
    & & & \multicolumn{2}{p{0.3\linewidth}}{
        \cellcolor[rgb]{0.9,0.9,0.9}{
            \makecell[{{p{\linewidth}}}]{
                \texttt{\tiny{[GM$|$GM]}}
                \texttt{Penalty: (2,4) is not empty, but contains 'D'.} \\
\\ 
\texttt{You have collectively accumulated 17 of 16 penalties.} \\
\texttt{Please try again!} \\
            }
        }
    }
    & & \\ \\

    \theutterance \stepcounter{utterance}  
    & & & \multicolumn{2}{p{0.3\linewidth}}{
        \cellcolor[rgb]{0.9,0.9,0.9}{
            \makecell[{{p{\linewidth}}}]{
                \texttt{\tiny{[GM$|$GM]}}
                \texttt{Maximum number of penalties exceeded} \\
            }
        }
    }
    & & \\ \\

    \theutterance \stepcounter{utterance}  
    & & & \multicolumn{2}{p{0.3\linewidth}}{
        \cellcolor[rgb]{0.9,0.9,0.9}{
            \makecell[{{p{\linewidth}}}]{
                \texttt{\tiny{[GM$|$GM]}}
                \texttt{* success: False} \\
\texttt{* lose: True} \\
\texttt{* aborted: True} \\
\texttt{{-}{-}{-}{-}{-}{-}{-}} \\
\texttt{* Shifts: 4.00} \\
\texttt{* Max Shifts: 12.00} \\
\texttt{* Min Shifts: 6.00} \\
\texttt{* End Distance Sum: 18.13} \\
\texttt{* Init Distance Sum: 30.19} \\
\texttt{* Expected Distance Sum: 29.33} \\
\texttt{* Penalties: 17.00} \\
\texttt{* Max Penalties: 16.00} \\
\texttt{* Rounds: 12.00} \\
\texttt{* Max Rounds: 28.00} \\
\texttt{* Object Count: 7.00} \\
            }
        }
    }
    & & \\ \\

    \theutterance \stepcounter{utterance}  
    & & & \multicolumn{2}{p{0.3\linewidth}}{
        \cellcolor[rgb]{0.9,0.9,0.9}{
            \makecell[{{p{\linewidth}}}]{
                \texttt{\tiny{[GM$|$GM]}}
                \texttt{Initial Grids:} \\
\texttt{```} \\
\texttt{Player 1:     Player 2:} \\
\texttt{ 1234567       1234567} \\
\texttt{╔═══╤═╤═╗     ╔═══╤═╤═╗} \\
\texttt{║◌◌◌│◌│◌║ 1   ║◌MS│◌│◌║ 1} \\
\texttt{║◌◌M│◌├─╢ 2   ║◌O◌│◌├─╢ 2} \\
\texttt{║AT◌├─┘◌║ 3   ║◌◌◌├─┘◌║ 3} \\
\texttt{║◌◌P│◌◌◌║ 4   ║◌◌◌│◌◌◌║ 4} \\
\texttt{╟─┬─┼───╢ 5   ╟─┬─┼───╢ 5} \\
\texttt{║◌│◌│S◌O║ 6   ║◌│◌│◌◌D║ 6} \\
\texttt{║◌│◌│◌D◌║ 7   ║T│◌│◌PA║ 7} \\
\texttt{╚═╧═╧═══╝     ╚═╧═╧═══╝} \\
\texttt{```} \\
\\ 
\texttt{Final Grids:} \\
\texttt{```} \\
\texttt{Player 1:     Player 2:} \\
\texttt{ 1234567       1234567} \\
\texttt{╔═══╤═╤═╗     ╔═══╤═╤═╗} \\
\texttt{║◌◌◌│◌│◌║ 1   ║◌◌◌│◌│◌║ 1} \\
\texttt{║◌◌M│◌├─╢ 2   ║◌O◌│◌├─╢ 2} \\
\texttt{║◌T◌├─┘◌║ 3   ║◌◌◌├─┘◌║ 3} \\
\texttt{║ADP│◌◌◌║ 4   ║ADM│◌◌◌║ 4} \\
\texttt{╟─┬─┼───╢ 5   ╟─┬─┼───╢ 5} \\
\texttt{║◌│◌│S◌O║ 6   ║◌│S│◌P◌║ 6} \\
\texttt{║◌│◌│◌◌◌║ 7   ║T│◌│◌◌◌║ 7} \\
\texttt{╚═╧═╧═══╝     ╚═╧═╧═══╝} \\
\texttt{```} \\
            }
        }
    }
    & & \\ \\

\end{supertabular}
}

\end{document}
